\chapter{材质、光照与逆渲染}

\paragraph*{本章阅读锚点}
原始 3DGS 的颜色通常由球谐系数直接表示，适合新视角合成，却把材质、光照、阴影、曝光和视角相关高光混在同一个外观参数里。只要用户想换灯光、编辑材质、处理镜子或 HDR，问题就会变成逆渲染：哪些是物体固有反照率，哪些是环境光，哪些是镜面反射，哪些只是相机曝光。本章按“把什么从颜色里拆出来”来读：材质分解方法拆 BRDF 和法线；反射方法拆镜面、虚像或高光；重光照方法拆光照传输；HDR 方法拆相机响应和曝光；物理路径追踪方法则把高斯重新解释成可追踪的几何和辐射基元。

\section{材质分解与逆渲染}

材质分解路线的起点不是“给 3DGS 多加几个颜色通道”，而是把原始 3DGS 的颜色混合公式重新解释清楚。普通 3DGS 在一个像素上的颜色可写成
\[
  \hat{\mathbf C}(\mathbf u)=\sum_i T_i(\mathbf u)\alpha_i(\mathbf u)\mathbf c_i(\mathbf u),
  \qquad
  T_i(\mathbf u)=\prod_{j<i}\left(1-\alpha_j(\mathbf u)\right).
\]
如果 $\mathbf c_i$ 只是球谐颜色，那么材质、环境光、阴影、曝光和高光都会被压进同一个外观项。逆渲染要做的是把 $\mathbf c_i$ 拆成几何、材质和光照，再用近似渲染方程重新合成：
\[
  L_o(x,\omega_o)=\int_{\Omega} f_r(x,\omega_i,\omega_o)L_i(x,\omega_i)
  \max(0,\mathbf n\cdot\omega_i)\,d\omega_i .
\]
这一步看似只是换颜色模型，实际会牵动三个基础问题：高斯有没有可用法线，前向 splatting 能不能获得遮挡/间接光信息，以及着色是在每个高斯上先算、还是先混合成像素级 G-buffer 后再算。

因此本节先读 2023.11 的两个早期锚点。GS-IR 把 3DGS 正式接入逆渲染，解决“怎么从高斯得到法线、材质和光照”的问题；GaussianShader 则从反射表面出发，说明“高斯级 shading function 怎样依赖稳定法线”。随后 DeferredGS、GI-GS、SVG-IR 等方法分别沿着延迟着色、全局光照、空间变化材质继续改造这条管线。



% METHOD: 2311.16473
\paragraph{GS-IR（3DGS 逆渲染锚点）（2023.11）}
GS-IR~\cite{2311_16473} 是把 3DGS 引入逆渲染的关键早期方法。原始 3DGS 能快速拟合新视角颜色，但它没有显式表面法线，也没有像 ray tracing 那样沿入射方向查询遮挡；如果直接给每个高斯加 BRDF 参数，优化会把阴影、反照率和间接光混在一起。GS-IR 的改变是把普通 3DGS 拆成三阶段：先得到可用几何和法线，再烘焙遮挡，最后用 PBR 管线优化材质与光照。

第一阶段仍从 calibrated multi-view RGB 输入训练 3D Gaussians，但高斯不只保存位置、尺度、旋转、不透明度和 SH 颜色。方法从渲染深度图的局部导数中构造法线监督，把深度梯度集中到真实表面附近，约束每个高斯携带的 normal 不要变成任意可学习向量。第二阶段处理 3DGS 的前向 splatting 缺口：NeRF/路径追踪可以沿 $\omega_i$ 反向采样遮挡，而 3DGS 只是把基元投到屏幕上，因此 GS-IR 预先把遮挡烘焙到球谐/体积缓存中，渲染时查询这个 occlusion cache 近似间接光和可见性。第三阶段关闭单纯颜色拟合，为高斯优化 albedo、roughness、metallic 等材质属性和环境光贴图，再按 PBR shading 合成观测颜色。

代码仓库也能看到这条三阶段流程：README 给出 Stage1 初始训练、\texttt{baking.py} 遮挡烘焙、Stage2 decomposition、\texttt{render.py}/\texttt{relight.py} PBR 渲染与重光照；源码中还有 \texttt{gs\_ir/volumes.py}、\texttt{occlusion\_kernel.cu}、\texttt{irradiance\_kernel.cu} 和 \texttt{pbr/shade.py} 等模块。这个实现结构说明 GS-IR 的贡献不只是“给 3DGS 加材质标签”，而是把法线估计、遮挡缓存、环境光和 BRDF 优化接成一条可复现的逆渲染流水线。它的局限也很清楚：遮挡和间接光是烘焙近似，镜面互反射和空间变化材质仍要靠后续方法补足。
\textbf{方法信息：} 官方代码：\url{https://github.com/lzhnb/GS-IR}；GitHub 约 439 stars；OpenAlex 约 84 次引用。



% METHOD: 2311.17977
\paragraph{GaussianShader（反射表面着色函数）（2023.11）}
GaussianShader~\cite{2311_17977} 解决的是另一条早期瓶颈：原始 3DGS 的 SH 颜色可以拟合中低频视角变化，但反射表面的高光位置由表面法线、观察方向和环境光共同决定；如果没有法线，模型常用错误几何或浮点高斯去“画”高光。GaussianShader 的改变是在每个高斯上估计法线，并用简化 shading function 替代纯 SH 颜色，使反射外观和表面朝向发生联系。

它保留 3DGS 的高效 rasterization 和 $\alpha$-混合，但把每个高斯的颜色改成由 diffuse color、roughness、specular tint、normal、residual color 和环境光共同决定。关键是法线参数化。训练后的 3D 高斯往往会贴着表面变扁，因此最短轴方向可以作为表面法线候选；GaussianShader 再学习一个 residual normal 修正最短轴误差，可概括为
\[
  \mathbf n_i=\mathrm{norm}\!\left(\mathbf v^{\min}_i+\Delta \mathbf n_i\right).
\]
为了避免 residual normal 随意吸收颜色误差，方法加入 normal-geometry consistency：把渲染深度图导出的像素法线作为几何参照，约束高斯法线和实际表面变化一致。渲染时每个高斯先根据法线、观察方向和可微环境光计算简化反射项，再把复杂反射中尚未解释的部分交给 residual color。

从公开代码看，\texttt{train.py} 的训练入口直接包含预测法线损失、残差法线损失和 zero-one 约束，命令行参数也单独控制法线权重、BRDF 环境和 BRDF 维度；渲染模块同时提供普通渲染和 lighting 渲染路径。这些实现细节对应论文中的设计：先让高斯获得稳定表面朝向，再把反射当作依赖法线的着色问题。相对 GS-IR，GaussianShader 更偏高斯级前向着色；相对后续 DeferredGS/GUS-IR，它的短板是多个高斯先各自着色再混合，像素级材质和法线可能被污染，但它为反射 3DGS 建立了“最短轴法线 + shading function”的重要锚点。
\textbf{实现细节：} 方法保留 3DGS 的显式高斯、tile rasterization 和快速 alpha 混合流程，但在颜色计算前为每个高斯构造表面法线。它利用训练中高斯趋于扁平的性质，把最短轴方向作为法线候选，并通过法线残差和几何一致性约束修正离散基元的法线不稳定。渲染时不再完全依赖 SH 拟合视角相关颜色，而是在高斯级别计算简化着色函数，再按 3DGS 的透明度顺序混合。这样能在保持训练和渲染效率的同时，把反射高光与表面朝向联系起来。 核心贡献是提出适合离散高斯基元的法线估计框架：利用 3D 高斯最短轴作为几何法线线索，并设计一致性损失约束法线与高斯几何关系，再将法线送入轻量 shading function 处理反射外观。 方法强调高斯级前向着色，强镜面表面中多个基元混合后的像素级法线、高光锐度和全局光照仍可能不足，后续 DeferredGS、GUS-IR、GlossyGS 等方法继续处理这些问题。
\textbf{实现补充：} 实现上还要核对输入、输出、参数、特征、初始化、训练损失、正则、渲染或推理接口，以及优化中的剪枝、分裂、采样、编码或解码步骤；这些抓手决定该方法到底改了 3DGS 管线的哪一环。
\textbf{方法信息：} 官方代码：\url{https://github.com/Asparagus15/GaussianShader}；GitHub 约 403 stars；OpenAlex 约 139 次引用。



% METHOD: 2404.09412
\paragraph{DeferredGS（延迟着色的解耦高斯）（2024.04）}
DeferredGS~\cite{2404_09412} 是本章最重要的早期锚点之一。它指出 3DGS 前向着色会产生混合伪影：每个高斯先根据自己的材质和光照着色，再按不透明度混合；重光照时，不同深度和材质的颜色已经被混在一起。DeferredGS 的改变是先渲染 G-buffer，再统一着色。

每个高斯不再只存 SH 颜色，而是存漫反射反照率、粗糙度和高光反照率等纹理参数。法线由联合训练的 SDF 网络蒸馏得到：SDF 的梯度提供连续法线先验，帮助离散高斯获得稳定表面方向。渲染时，系统先得到位置、法线和材质图，再在像素空间用可学习环境贴图执行着色，支持纹理编辑和重光照。后续 Ref-Gaussian、RGS-DR、MaterialRefGS 等都可看作在这个延迟着色思路上强化反射、间接光或材质一致性。
\textbf{实现细节：} DeferredGS（2024.04） （Decoupled and Editable）将延迟着色引入 3DGS 材质分解。直接在逐高斯上进行前向着色（forward shading）再 $\alpha$-混合 的方法存在致命缺陷：即使法线蒸馏损失在像素级约束了混合法线的正确性，每个高斯的个体法线和着色颜色可能仍然错误，因为几何属性（位置、协方差、不透明度）可能过拟合到特定光照条件以产生正确的混合法线——换光照后混合就失效。DeferredGS 的解决方案是延迟着色：先将位置 $P$、法线 $n$、漫反射反照率 $k_d$、粗糙度 $r$、镜面反射率 $k_s$ 和不透明度 $\alpha$ 通过 $\alpha$-混合 投影为 2D 像素属性图 $I_P, I_n, I_{k_d}, I_r, I_{k_s}, I_\alpha$，然后对 $I_\alpha > 0.5$ 的像素用 Disney BRDF 着色。光照用 $6\times512\times512$ HDR cubemap 建模，漫反射辐照 $c_{diff} = k_d \int L_i(\omega_i)(\omega_i \cdot n) d\omega_i$ 预计算为 2D 纹理，镜面反射采用 split-sum 近似分解为预过滤环境光贴图（mipmap，每级对应固定粗糙度）和 BRDF 积分 LUT。法线估计依赖联合训练的 SDF 网络（基于 multi-level hash encoding 的 NeuS 架构）：SDF 的体渲染密度 $\rho(t) = \max(-\frac{d}{dt}\Phi_s(f(x(t))), 0)$（$\Phi_s$ 为 Sigmoid），法线蒸馏损失 $L_{nd} = \sum_v \sum_{i=1}^M \|1 - C_n(v) \cdot \nabla f_s(x_{v,i})\|$ 将 SDF 法线场的连续性注入离散高斯。纹理编辑通过从输入视角绘制修改后的纹理图，然后用 ASR（As-Rigid-As-Possible）变形代理 mesh 扩展到随机视角优化。
\textbf{设计动机与边界：} 相对原始 3DGS 以独立高斯和图像重建误差为主的设定，DeferredGS 把问题落在以下可实现的改动上：DeferredGS方法包含四个核心模块：（1）纹理定义：在每个3D高斯上定义可优化的纹理参数（漫反射反bedo kd、粗糙度r、高光反照率ks）替代原始SH，使纹理和光照可分离。（2）法线场蒸馏：直接从离散高斯提取法线质量差（尤其在反射面上）。联合训练一个基于多级哈希编码的SDF网络，利用SDF的梯度作为法线先验引导高斯的法线优化，从连续隐式场蒸馏出更优的表面法线。同时分离SDF颜色分支为视角无关和视角相关两部分，避免几何歧义。（3）延迟着色：不同于前向着色先对每个高斯着色再混合（在重光照时因不透明度在新光照下不正确导致混合伪影），DeferredGS先渲染G-buffer（法线、位置、纹理参数），再在像素空间统一着色，避免混合伪影。（4）用可学习环境贴图建模场景光照。支持重光照和纹理编辑应用。其设计目的不是单独增加一个模块，而是用上述表示、监督或渲染约束消除原流程在目标任务中的具体失配。提出DeferredGS，通过延迟着色解决3DGS纹理-光照解耦中的混合伪影问题。核心创新包括：（1）联合训练SDF网络进行法线场蒸馏，从连续隐式场引导离散高斯获得更优表面法线；（2）在高斯上绑定纹理参数（漫反射反照率、粗糙度、高光反照率）替代SH；（3）采用延迟着色先渲染后着色，避免前向着色的混合伪影。实现高质量的纹理编辑和重光照效果。适用时还应保留论文给出的边界：SDF网络增加训练开销；对多光源场景的环境贴图建模能力有限。
\textbf{方法信息：} ByteDance 参与了这篇论文；作者关联代码：\url{https://github.com/IGLICT/DeferredGS}。



% METHOD: 2408.03538
\paragraph{PRTGS（预计算辐射传输重光照）（2024.08）\cite{2408_03538}}
PRTGS 把传统 Precomputed Radiance Transfer 引入 3DGS，解决动态环境光下软阴影和间接光难以实时计算的问题。GS-IR 和 Relightable 3D Gaussian 一类方法可以估计材质和环境光，但当光照改变时，如果每帧重新 ray tracing 或重新估计遮挡，速度很难保持实时；如果只用环境贴图直接着色，阴影和互反射又会过于粗糙。PRTGS 的改变是把昂贵的光传输提前算好，并把每个高斯的传输函数压到球谐系数中。

每个高斯携带 normal、visibility 和 Disney BRDF 参数。训练或预计算阶段，系统为高斯计算阴影、遮挡和互反射相关的 transfer：漫反射表面可表示为 SH 向量，光泽表面则使用更高维的 transfer matrix。渲染阶段，当环境光变化时，不再逐像素追踪大量二次光线，而是把环境光 SH 系数与每个高斯的 transfer 向量/矩阵做点积，得到动态光照下的出射辐射，再进入 3DGS splatting。由于光传输近似满足线性叠加，这种表示可以在 1080p 下达到 30 fps 级动态重光照。相对 GI-GS 的训练时路径追踪，PRTGS 更适合交互式换光；代价是主要面向低频环境光，且预计算阶段需要为当前几何存储额外传输数据。当前只确认论文材料，未确认完整公开代码。

\textbf{实现细节：} PRTGS（2024.08） （Precomputed Radiance Transfer）修改了原始 3DGS 的着色环节——用基于 Disney BRDF 的完整 PRT 管线替代纯漫反射假设——高斯的几何与位置属性保持不变。传统做法（如 GS-IR、RelightableGS）假设高斯表面为纯漫反射，无法处理镜面互反射。PRTGS 的核心改进是为每个高斯建立基于 Disney BRDF 的完整着色模型，将 BRDF 分解为漫反射项 $f_d = (1-m)/(4\pi)$ 和镜面微面项 $f_s$（含法线分布 $D$、菲涅尔 $F$、几何遮蔽 $G$），并为每个高斯优化反照率、粗糙度 $r \in [0,1]$ 和金属度 $m \in [0,1]$ 三个材质参数。自传输递归。预计算阶段的关键创新是自传输递归：不是直接计算间接光照（计算量随弹射次数指数增长），而是对每个高斯执行一次弹射的光线追踪，构建索引矩阵记录每个方向上贡献最大的遮挡高斯，然后递归计算 $n$-弹射自传输向量 $\mathbf{t} = \sum_{j=1}^{n} \mathbf{t}^j$，其中一弹射项 将漫反射传输向量逐弹射累积。
\textbf{方法信息：} 当前未确认企业背景、完整代码、高星标或高引用信息。入选依据：企业:否；代码:否；星标:否；引用:否。
\textbf{设计动机与边界：} 相对原始 3DGS 以独立高斯和图像重建误差为主的设定，PRTGS 把问题落在以下可实现的改动上：方法为每个3D高斯溅射分配法线、可见性和BRDF属性。通过深度梯度推导伪法线作为优化引导。预计算阴影和相互反射等复杂传输函数所需的昂贵传输模拟，结果以球谐(SH)向量或矩阵表示。提出针对3DGS的独特射线追踪和间接光照预计算技术：训练阶段使用基于射线的精确可见性计算，渲染阶段使用预计算的传输函数进行快速着色。自传输递归计算间接光照，将辐射度传输投影到SH系数实现快速渲染。采用Disney BRDF模型，包含albedo、roughness和metallic参数。其设计目的不是单独增加一个模块，而是用上述表示、监督或渲染约束消除原流程在目标任务中的具体失配。首次将预计算辐射度传输(PRT)引入3DGS，通过预计算每个高斯溅射的辐射度传输函数（含阴影和相互反射），在低频光照环境下实现30+fps实时高质量重新光照，支持1080p分辨率下相对复杂的场景。适用时还应保留论文给出的边界：适用于低频光照环境，高频光照效果有限；预计算阶段耗时；球谐阶数限制高频细节表达
% METHOD: 2408.03753
\paragraph{3iGS（张量分解入射照明）（2024.08）}
3iGS~\cite{2408_03753} 改的是原始 3DGS 的视角相关颜色模型。逐高斯独立球谐系数难以表达连续照明场和 BRDF 共同作用下的镜面变化；3iGS 把出射辐射度分解成漫反射色和视角相关镜面色，并显式查询局部入射照明场。

实现上，3iGS 为每个高斯保留位置、opacity 和几何参数，但把原本直接学习的 SH 颜色换成三部分：稳定的 diffuse color、每高斯 BRDF feature，以及一个连续的局部入射照明场。照明场借鉴 TensoRF 的向量-矩阵分解，用较粗的三维张量网格表示；训练时用高斯中心坐标查询该网格，得到该空间位置的 incident illumination feature。随后一个小型神经渲染器接收 BRDF feature、入射照明和观察方向，输出 specular color，再与 diffuse color 合成该高斯的颜色并进入普通 3DGS rasterization。相对 DeferredGS 的 PBR G-buffer，3iGS 不要求把材质分解成可解释的粗糙度或金属度，而是用“局部照明场 + BRDF 特征”提升复杂视角相关外观。
\textbf{方法信息：} 代码：\url{https://github.com/tangzj/3igs}。
\textbf{设计动机与边界：} 相对原始 3DGS 以独立高斯和图像重建误差为主的设定，3iGS 把问题落在以下可实现的改动上：核心思想是将高斯出射辐射度分解为漫反射色和视角相关镜面色两部分。对于每个高斯，入射照明场通过连续网格表示的因式分解张量计算，借鉴TensoRF的向量-矩阵空间分解：每个照明张量由R个分量组成，每个分量包含特征向量和特征矩阵的组合。高斯的中心坐标作为输入查询因式分解张量网格获取照明特征。然后通过小型神经网络F将BRDF特征、入射照明和观察角度映射为镜面色。照明场作为连续网格函数联合优化。其设计目的不是单独增加一个模块，而是用上述表示、监督或渲染约束消除原流程在目标任务中的具体失配。提出3iGS方法，将3DGS中每个高斯的出射辐射度表达为局部照明场和BRDF特征的函数，通过类似TensoRF的向量-矩阵空间分解学习连续入射照明场，显著增强3DGS的镜面视角相关渲染效果，同时保持快速训练和渲染速度。适用时还应保留论文给出的边界：不直接分解为物理BRDF属性；对极端镜面反射效果仍有局限；张量分解分辨率可能限制高频照明细节
% METHOD: 2409.10335
\paragraph{Phys3DGS（隐藏高斯与混合网格-高斯逆渲染）（2024.09）\cite{2409_10335}}
Phys3DGS 接着 DeferredGS 的问题继续往下拆。GS-IR、GaussianShader 这类方法通常仍沿用 3DGS 的体式 $\alpha$ 混合：同一条射线上的多个高斯都会参与颜色累积。做 NVS 时这有利于拟合外观，但做 PBR 重光照时，表面下方或错误深度处的 hidden Gaussians 也会带着法线、BRDF 和可见性参与着色，导致换光后阴影和高光被污染。Phys3DGS 的改变是明确把这个问题归因到渲染顺序，并把逆渲染改成“先确定表面，再在表面上着色”的延迟路线。

具体流程分两步。第一步仍训练 3DGS 并从中提取 mesh；mesh 不是最终渲染资产，而是给高斯提供更稳定的表面几何、法线和可见性参照。第二步使用 hybrid mesh-3DGS representation：高斯保留快速 splatting 的外观和可编辑材质参数，mesh 则参与约束法线、深度和可见性，避免隐藏高斯在重光照时提供错误贡献。渲染时不再让每个高斯先独立求 BRDF 后混合，而是先得到像素级几何/材质缓冲，再按 PBR 方程求颜色。相对 DeferredGS，Phys3DGS 的重点更具体：它把“延迟着色为什么必要”落实到 hidden Gaussians 对 relighting 的破坏，并用 mesh 提供几何校准。当前只确认论文材料，未确认完整公开代码。

\textbf{实现细节：} Phys3DGS（2024.09） （Physically-based 3D Gaussian Splatting）修改了原始 3DGS 的外观表示环节——用物理 PBR 参数（反照率、粗糙度、金属度）替代球谐颜色编码——几何重建和泼溅渲染管线保持不变。该方法采用两阶段训练策略，逐步从"重建外观"过渡到"分解材质"。第一阶段：几何与混合表示。先用标准 3DGS 训练并提取网格，得到 mesh-3DGS 混合表示：每个三角面片上分配一个扁平高斯，位置和形状由面片顶点决定。为解决 SDF 正则化导致的大高斯问题，提出 masked opacity 正则化 $L_o = -(1 - M) \log(1 - o)$（其中 $M$ 为二值掩码）。直觉上，这个损失只约束"应该透明"的背景区域，对前景高斯"放手不管"，从而不干扰后续 BRDF 参数的学习。同时引入高斯尺度正则化 $L_{sc} = \sum_i \max((s^i) - \lambda_{r_c} r^i_{sc}, 0)$，其中 $r^i_{sc}$ 为面片外接圆半径，$\lambda_{r_c} = 0.2$，防止高斯超出其对应三角面片的范围。
\textbf{方法信息：} 当前未确认企业背景、完整代码、高星标或高引用信息。入选依据：企业:否；代码:否；星标:否；引用:否。


% METHOD: 2410.02619
\paragraph{GI-GS（全局光照分解）（2024.10）}
GI-GS~\cite{2410_02619} 针对大多数逆渲染方法忽略间接光照的问题。只建模直接环境光时，墙面反弹光、物体间遮挡和颜色 bleeding 会被误分到反照率或残差里。GI-GS 的改变是在 3DGS 逆渲染中加入高效路径追踪计算间接光。

第一阶段用 vanilla 3DGS 重建几何，并为高斯增加法线属性，用深度伪法线和 TV 损失平滑法线。第二阶段用延迟着色渲染深度、法线和 BRDF 图，基于 IBL 和可学习环境图分解直接漫反射/镜面反射。第三阶段固定已有几何和直接光照解释，用轻量可微路径追踪从 G-buffer 出发估计一次或多次反弹的间接光，再把直接项和间接项合成训练图像。这样做的好处是，换环境光时 indirect illumination 会随新光照重新计算，而不是仍使用训练时烘焙在高斯颜色里的反弹光。相对 GS$^3$ 的三重 splatting 近似，GI-GS 更接近物理全局光照；相对 DeferredGS，它把“只在像素级做一次 PBR”推进到“用 path tracing 修正间接光”。
\textbf{实现细节：} GI-GS（2024.10） （Global Illumination Decomposition）通过三阶段管线将全局光照显式分解为直接光照、阴影和间接光照。Baseline 方法（如 GS-IR）用球谐编码间接光照，但球谐缺乏物理约束，学到的间接光照不可重光照。GI-GS 的设计动机是：将每个物理成分建模为独立的可微操作，使分解结果可编辑、可重光照。第一阶段用标准3DGS重建几何，法线作为高斯属性优化，通过深度图导出的伪法线 $\hat{n}$ 监督渲染法线：$L_n = \|n - \hat{n}\| + \lambda_{TV} L_{TV_{\mathrm{normal}}}$，避免强制高斯退化为扁平盘的几何正则化导致的渲染质量下降。第二阶段建模直接光照：镜面分量通过 split-sum 近似分解为 $L_s = \underbrace{R}_{\mathrm{BRDF积分}} \cdot \underbrace{I_s}_{\mathrm{预滤波环境贴图}}$，漫反射分量从预卷积环境贴图 $I_d$ 查询。环境光遮蔽 $O(\mathbf{p})$ 通过从表面点向半球均匀采样射线、射线行进检测交点来计算：$O(\mathbf{p}) = 1 - \frac{1}{\pi}\int_{\Omega^{}+} V(\mathbf{p}, \omega) (\omega \cdot n) d\omega$，其中可见性 $V=1$ 当射线与表面相交。直接光照最终为 $L_{\mathrm{dir}} = (1-m)\boldsymbol{a} \cdot O(\mathbf{p}) \cdot I_d + R \cdot I_s$。第三阶段建模间接光照：从第一阶段渲染的直接光照图像 $I_{\mathrm{dir}}$ 出发，对每个表面点做路径追踪，命中点的出射辐射度直接索引 $I_{\mathrm{dir}}$ 中对应像素值：$L_i(\mathbf{p}, \omega_i) = V(\mathbf{p}, \omega_i) \cdot I_{\mathrm{dir}}(\mathbf{p}')$。最终渲染为两阶段之和 $L_o = L_{\mathrm{dir}} + (1-m)\boldsymbol{a} \cdot L_{\mathrm{ind}}$。对于复杂场景，GI-GS 进一步将屏幕空间路径追踪扩展到基于立方体图（Cubemap）的世界空间，渲染六面深度图和 RGB 图覆盖完整场景几何，使间接光照能捕获视锥外的遮挡和反射。
\textbf{设计动机与边界：} 相对原始 3DGS 以独立高斯和图像重建误差为主的设定，GI-GS 把问题落在以下可实现的改动上：GI-GS分三阶段：(1)几何重建：使用vanilla 3DGS重建场景几何，将法线作为每个高斯的新属性，通过深度图导出的伪法线进行监督，结合TV损失实现平滑法线；(2)直接光照建模：通过延迟着色渲染G-buffer（深度图、法线图、BRDF图），使用基于图像的光照(IBL)和可学习环境图进行PBR渲染，分解漫反射和镜面反射分量；(3)间接光照计算：基于几何和直接光照结果，使用可微分路径追踪计算间接光照，结合直接和间接渲染结果进行最终优化。其设计目的不是单独增加一个模块，而是用上述表示、监督或渲染约束消除原流程在目标任务中的具体失配。提出GI-GS逆渲染框架，首次在3DGS中通过高效路径追踪和延迟着色实现准确的间接光照计算，使重光照时能正确建模复杂光物交互，在渲染质量和效率上均优于现有方法。适用时还应保留论文给出的边界：路径追踪增加渲染开销；依赖几何重建质量；复杂场景中的多次反弹间接光照可能计算昂贵
\textbf{方法信息：} 代码：\url{https://github.com/stopaimme/GI-GS}；GitHub 关注度较高，约 131 stars；引用很多，Semantic Scholar 约 43 次。



% METHOD: 2410.11419
\paragraph{GS$^3$（三重 splatting 重光照）（2024.10）}
GS$^3$~\cite{2410_11419} 处理点光源输入下的复杂外观，包括半透明、蓬松和各向异性材质。它不把重光照只写成一次材质着色，而是设计三重 Gaussian splatting 流程。

第一重 splatting 按正常视角渲染空间高斯，用于颜色计算；第二重把所有空间高斯朝光源方向投影，估计从点光源看过去的可见性，也就是自阴影，再用小型 MLP 修正阴影边界和半透明区域；第三重为每个空间高斯预测 RGB 补偿项，吸收多次反弹、半透明散射和其他未显式建模的残差。表示上，每个空间高斯不再只携带颜色，而是使用 Lambertian 项加角度 Gaussian mixture 作为反射函数，入射方向和观察方向一起决定输出。相对 GI-GS 的路径追踪间接光照，GS$^3$ 更偏实时重光照中的高效近似：牺牲一部分物理可解释性，换来点光源采集下的高速训练和 90 fps 级渲染。
\textbf{实现细节：} GS$^{}3$（2024.10） （Triple Gaussian Splatting）与前面方法的"统一 BRDF"路线不同，它将反射显式分解为四个独立的物理分量——漫反射、镜面反射、阴影和次表面散射（SSS）——每个分量有独立的学习基色和反射率强度。标准 3DGS 用球谐系数（SH）编码外观，无法将材质与光照解耦，因此不支持重光照；GS$^{}3$ 用这种分量化的物理着色模型替换 SH。具体地，每个高斯的颜色计算为 $C_i = (c_d f_d + c_s f_s) S(x) + c_{\mathrm{sss}} f_{\mathrm{sss}}$，其中 $c_d, c_s, c_{\mathrm{sss}} \in \mathbb{R}^{}3$ 分别为漫反射、镜面反射和 SSS 的学习基色，$f_d, f_s, f_{\mathrm{sss}}$ 为对应反射率标量，$S(x)$ 为阴影因子。漫反射项采用 Lambertian BRDF（视角无关的余弦加权反射）；镜面反射项用 Fresnel 因子调制各向异性球面高斯（ASG）基，能表达拉丝金属等各向异性高光。SSS 项采用偶极子扩散剖面（dipole diffusion），由一个 6 层 MLP（隐藏维度 256）从位置编码（$L=4$ 频率段）的坐标和方向预测散射系数 $\sigma_s \in [0.05, 2.05]$、吸收系数 $\sigma_a \in [0.05, 2.05]$ 和表面分离距离 $r \in [0.1, 3.1]$，再代入经典偶极子公式 $f_{\mathrm{sss}}(r) = \frac{\sigma_t'}{4\pi}\!\left(\frac{z_r(\sigma_t d_r + 1)e^{-\sigma_t d_r}}{d_r^{}3} + \frac{z_r z_v(\sigma_t d_v + 1)e^{-\sigma_t d_v}}{d_v^{}3}\right)$ 计算散射强度，其中 $\sigma_t' = \sigma_s + \sigma_a$，$z_r, z_v$ 为实偶极子和虚偶极子源深度。阴影采用两阶段方法：先沿光线累积高斯不透明度得到粗糙透射率 $v_i = \prod_{k \in O_i}(1-\alpha_k)$，再以投影密度加权聚合为 $v_g$，最后用 3 层 MLP（32 维隐藏层，位置编码 $L=3$）精修为 $S(x)$。
\textbf{锚点细化：} GS$^{}3$：（Triple Gaussian Splatting）与前面方法的"统一 BRDF"路线不同，它将反射显式分解为四个独立的物理分量——漫反射、镜面反射、阴影和次表面散射（SSS）——每个分量有独立的学习基色和反射率强度。标准 3DGS 用球谐系数（SH）编码外观，无法将材质与光照解耦，因此不支持重光照；GS$^{}3$ 用这种分量化的物理着色模型替换 SH。具体地，每个高斯的颜色计算为 $C_i = (c_d f_d + c_s f_s) S(x) + c_{\mathrm{sss}} f_{\mathrm{sss}}$，其中 $c_d, c_s, c_{\mathrm{sss}} \in \mathbb{R}^{}3$ 分别为漫反射、镜面反射和 SSS 的学习基色，$f_d, f_s, f_{\mathrm{sss}}$ 为对应反射率标量，$S(x)$ 为阴影因子。漫反射项采用 Lambertian BRDF（视角无关的余弦加权反射）；镜面反射项用 Fresnel 因子调制各向异性球面高斯（ASG）基，能表达拉丝金属等各向异性高光。SSS 项采用偶极子扩散剖面（dipole diffusion），由一个 6 层 MLP（隐藏维度 256）从位置编码（$L=4$ 频率段）的坐标和方向预测散射系数 $\sigma_s \in [0.05, 2.05]$、吸收系数 $\sigma_a \in [0.05, 2.05]$ 和表面分离距离 $r \in [0.1, 3.1]$，再代入经典偶极子公式 $f_{\mathrm{sss}}(r) = \frac{\sigma_t'}{4\pi}\!\left(\frac{z_r(\sigma_t d_r + 1)e^{-\sigma_t d_r}}{d_r^{}3} + \frac{z_r z_v(\sigma_t d_v + 1)e^{-\sigma_t d_v}}{d_v^{}3}\right)$ 计算散射强度，其中 $\sigma_t' = \sigma_s + \sigma_a$，$z_r, z_v$ 为实偶极子和虚偶极子源深度。阴影采用两阶段方法：先沿光线累积高斯不透明度得到粗糙透射率 $v_i = \prod_{k \in O_i}(1-\alpha_k)$，再以投影密度加权聚合为 $v_g$，最后用 3 层 MLP（32 维隐藏层，位置编码 $L=3$）精修为 $S(x)$。训练采用渐进式策略，按漫反射$\to$阴影$\to$SSS$\to$镜面反射的顺序逐步激活各分量。这种分量分离使重光照只需改变光照输入，材质参数保持不变。
\textbf{设计动机与边界：} 相对原始 3DGS 以独立高斯和图像重建误差为主的设定，GS\^{}3 把问题落在以下可实现的改动上：方法包含：（1）表示：每个空间高斯使用Lambertian+角度高斯混合作为反射函数，描述复杂外观。（2）三重splatting过程：第一重splatting渲染正常视角图像用于颜色计算；第二重splatting将所有空间高斯投影向光源方向获取阴影值，通过小型MLP精炼自阴影效果；第三重splatting通过网络计算每个空间高斯的RGB补偿元组，补偿全局光照等效果。（3）训练和推理：在单张消费级GPU上40-70分钟训练，90fps渲染。其设计目的不是单独增加一个模块，而是用上述表示、监督或渲染约束消除原流程在目标任务中的具体失配。提出GS\^{}3三重高斯splatting表示和渲染方法，实现90fps实时高质量重光照，支持从固体到蓬松、从半透明到各向异性的广泛外观变化，训练时间40-70分钟。适用时还应保留论文给出的边界：对高度复杂的全局光照效果（如焦散）建模有限；点光源假设限制了场景光照类型
\textbf{方法信息：} 代码：\url{https://github.com/gsrelight/gs-relight}；GitHub 关注度较高，约 178 stars；引用很多，Semantic Scholar 约 56 次。



% METHOD: 2410.13349
\paragraph{GlossyGS（光泽物体的微面先验）（2024.10）}
GlossyGS~\cite{2410_13349} 针对光泽物体的几何-材质歧义。镜面反射会把环境内容映射到物体表面，优化时可能误以为是几何细节或纹理。GlossyGS 的改变是加入微面几何分割先验，并重新安排法线和材质的着色顺序。

方法先混合法线图，再生成材质图，最后执行 Cook-Torrance BRDF 着色，以更准确模拟反射面法线分布。微面几何分割先验由视觉基础模型的语义表示训练而来，用于判断哪些区域应当共享相近的微面法线统计，避免把环境反射纹理误当成真实几何凹凸。表示上，它结合显式 anchor 的宏观几何和隐式特征预测的多个 neural Gaussians：anchor 给出粗表面位置，关联特征再解码出局部神经高斯和 BRDF 参数，用来表达微面尺度的法线/材质变化。相对 Spec-Gloss Surfels 的 BRDF 参数化选择，GlossyGS 更重视法线分布和语义微面先验；相对 Ref-Gaussian，它不以 inter-reflection 为主，而是优先减轻光泽物体的几何-材质歧义。
\textbf{方法信息：} Alibaba 参与了这篇论文；当前未确认完整公开代码。
\textbf{设计动机与边界：} 相对原始 3DGS 以独立高斯和图像重建误差为主的设定，GlossyGS 把问题落在以下可实现的改动上：方法包含三个核心组件：(1)法线图预滤波策略：修改着色顺序，先混合法线再生成材质图，最后着色渲染，解决了3DGS中高斯体着色顺序对光泽物体几何重建的负面影响；(2)微面几何分割先验：基于大视觉模型提取的语义表示训练分割模型，调整微面法线分布和反射面几何，显著减少逆向渲染的固有歧义；(3)混合几何与材质表示：使用锚点(anchor)显式表示宏观粗几何，通过关联的特征向量隐式预测多个神经高斯体(neural Gaussians)，精确模拟微面法线分布和材质属性。整体基于Cook-Torrance BRDF模型进行材质分解。其设计目的不是单独增加一个模块，而是用上述表示、监督或渲染约束消除原流程在目标任务中的具体失配。提出微面几何分割先验减少逆向渲染的固有歧义性，设计法线图预滤波策略准确模拟反射面法线分布，结合混合显式-隐式表示精确重建光泽物体的几何与材质属性适用时还应保留论文给出的边界：依赖自定义构建的数据集训练微面几何分割模型；对极端光泽/镜面物体（如完美镜面）的重建仍有挑战；逆向渲染的固有歧义性只能减轻无法完全消除
% METHOD: 2411.07478
\paragraph{GUS-IR（统一前向-延迟着色逆渲染）（2024.11）\cite{2411_07478}}
GUS-IR 试图调和两类着色路线。前向着色在每个高斯上直接计算 BRDF，能保留基元属性和一定局部表达，但多个高斯混合后容易把材质和法线糊在一起；延迟着色先 rasterize 出 G-buffer，再在像素级求 PBR，材质边界更清楚，却可能丢掉高斯级视角相关细节。GUS-IR 的改变是提出 unified shading：让几何、材质和光照分解同时利用前向着色和延迟着色的优势。

实现上，它继续使用最短轴作为高斯法线候选，并配合深度相关正则稳定形状；光照侧继承 GS-IR 的 probe/baking 思路，但强化环境遮挡，使粗糙和光泽表面上的间接光不会全部塞进反照率。统一着色可以理解为两条通道协同：高斯级分支保留与 primitive 绑定的材质和局部外观，像素级分支在混合后计算更一致的 BRDF 响应。相对 GaussianShader，GUS-IR 不只给每个高斯一个 shading function，而是显式讨论高斯先着色与像素后着色的取舍；相对 Phys3DGS，它更关注 shading strategy 本身，而不是用 mesh 解决 hidden Gaussians。当前只确认论文和项目材料，未确认完整公开代码。

\textbf{实现细节：} 前面的方法各自选择了前向着色或延迟着色中的一种。GUS-IR（2024.11） （Gaussian Splatting with Unified Shading）则认识到两者各有优劣，核心思想是：前向着色与延迟着色并非对立选择，而是分别适配漫反射与光泽表面的不同物理先验——统一框架的本质是通过可学习权重实现先验的自适应选择。核心创新是统一着色：同时计算前向和延迟两条着色路径，通过可学习的权重平衡两者。前向着色（如 GaussianShader）先在每个高斯上独立计算 BRDF 再 $\alpha$-混合——这有利于漫反射表面的反照率恢复，因为各基元独立着色后混合的颜色更接近真实漫反射。延迟着色（如 GS-IR）先将材质属性混合到像素级再计算 BRDF——这有利于光泽表面，因为混合后的单一法线能保持镜面高光的锐利。
\textbf{实现补充：} 实现上还要核对输入、输出、参数、特征、初始化、训练损失、正则、渲染或推理接口，以及优化中的剪枝、分裂、采样、编码或解码步骤；这些抓手决定该方法到底改了 3DGS 管线的哪一环。
\textbf{方法信息：} 当前未确认企业背景、完整代码、高星标或高引用信息。入选依据：企业:否；代码:否；星标:否；引用:否。
\textbf{设计动机与边界：} 相对原始 3DGS 以独立高斯和图像重建误差为主的设定，GUS-IR 把问题落在以下可实现的改动上：GUS-IR首先分析比较了前向着色和延迟着色在处理复杂材质时的优劣，然后提出统一着色方案融合两者优势实现更好的分解。在几何表示方面，分析3DGS中的法线建模，利用每个高斯基元的最短轴作为法线，配合深度相关正则化改进形状重建。在光照建模方面，增强GS-IR提出的基于探针的烘焙方案，实现更准确的环境光遮蔽建模以更好地处理间接光照。其设计目的不是单独增加一个模块，而是用上述表示、监督或渲染约束消除原流程在目标任务中的具体失配。提出统一前向-延迟着色方案用于3DGS逆渲染，利用高斯最短轴作为法线并引入深度正则化改进几何表示，增强探针烘焙方案实现更精确的环境光遮蔽建模。适用时还应保留论文给出的边界：论文Introduction部分被模板文本覆盖，具体实验数据受限
% METHOD: 2412.15867
\paragraph{IRGS（2D 高斯光线追踪与互反射）（2024.12）\cite{2412_15867}}
IRGS 处理的是 GS-IR 留下的另一个硬缺口：前向 splatting 很快，但它不能像 ray tracing 那样沿任意入射方向查询可见性和间接辐射。因此许多 3DGS 逆渲染方法只能简化渲染方程，或用可学习参数近似 indirect light。IRGS 的改变是把 2D Gaussian Splatting 改造成可微 ray tracing primitive，用 2D Gaussian ray tracing 查询 incident radiance，从而在高斯框架中计算更完整的 rendering equation。

管线先训练一个 2DGS 几何模型，得到稳定的深度和法线；随后给每个 2D Gaussian 分配 diffuse albedo 和 roughness，rasterization 得到像素级 albedo/roughness/normal。关键模块是 2DGRT：射线不再与体式 3D Gaussian 求一个模糊的最大响应点，而是与 2D Gaussian disk 求有定义的交点，并用透明度累积获得沿入射方向看到的颜色和可见性。渲染方程中，入射光被拆成来自环境贴图的直接项和由 2DGRT 查询的间接项：
\[
  L_i(\omega_i,x)=L_{\mathrm{dir}}(\omega_i)V(\omega_i,x)+L_{\mathrm{ind}}(\omega_i,x).
\]
训练时只对一部分像素采样多条入射光线，降低 Monte Carlo 开销；重光照时则用 ray tracing 查询到的 albedo、roughness 和 normal 估计间接辐射，避免完全忽略互反射。相对 GS-IR 的遮挡烘焙，IRGS 更直接地把可见性和间接光查询放进优化；相对 EAG-PT，它仍服务于逆渲染和重光照，而不是完整多次反弹路径追踪。官方代码公开在 \url{https://github.com/fudan-zvg/IRGS}，GitHub 约 166 stars；该方法目前是互反射高斯逆渲染路线中必须单独说明的锚点。

\textbf{实现细节：} IRGS（2024.12） （Inter-Reflective Gaussian Splatting）显式处理表面间的相互反射。传统 3DGS 逆渲染方法（如 R3DG）无法高效查询任意方向的可见性和间接辐射，只能采用 split-sum 近似或可学习参数建模间接光照。IRGS 提出在 2D 高斯基元上直接进行光线追踪（2D Gaussian Ray Tracing, 2DGRT），利用 2D 高斯盘具有明确定义的射线-面片交点这一特性，消除 3D 高斯交点不一致的问题。具体地，给定 2D 高斯的中心 $\mu$、两个切向量 $t_u, t_v$ 和缩放 $s=(s_u,s_v)$，射线-面片交点解析计算为 $p = r_o + t r_d$，其中 $t = (\mu - r_o) \cdot n / (r_d \cdot n)$，法线 $n = t_u \times t_v$。通过为每个 2D 高斯构建自适应二十面体包围体并建立 BVH，利用 NVIDIA OptiX 的硬件加速光线-三角形求交实现高效追踪。框架分两阶段训练：Stage I 用标准 2DGS 预训练获取可靠几何（含法线一致性损失 $L_n$、深度畸变损失 $L_d$、边缘感知法线平滑 $L_{s,n}$ 和 mask 熵损失 $L_o$）；Stage II 赋予每个高斯漫反射反照率 $a$ 和粗糙度 $r$ 参数，通过延迟着色在像素级评估完整渲染方程。入射辐射分解为直接辐射和间接辐射：$L_i(\omega_i, x) = V(\omega_i, x) L_{dir}(\omega_i) + L_{ind}(\omega_i, x)$，其中环境 cubemap 建模直接光照，可见性 $V$ 和间接辐射 $L_{ind}$ 通过 2DGRT 实时查询。每个像素用 Monte Carlo 采样 $N_r$ 个入射方向均匀采样半球，最终颜色 $c_{pbr} = \frac{2\pi}{N_r}\sum_{i=1}^{N_r}(f_d+f_s) L_i(\omega_i, x)(\omega_i \cdot n)$。
\textbf{实现补充：} 实现上还要核对输入、输出、参数、特征、初始化、训练损失、正则、渲染或推理接口，以及优化中的剪枝、分裂、采样、编码或解码步骤；这些抓手决定该方法到底改了 3DGS 管线的哪一环。
\textbf{方法信息：} 代码已公开，仓库为 fudan-zvg/irgs (166 stars)；GitHub 关注度很高，fudan-zvg/irgs 约 166 stars。入选依据：企业:否；代码:是；星标:是；引用:否。


% METHOD: 2504.06815
\paragraph{SVG-IR（空间变化材质逆渲染）（2025.04）}
SVG-IR~\cite{2504_06815} 的问题要放在 GS-IR 和 GaussianShader 之后看：这两个早期锚点通常把一个高斯看成一个材质常量单元，同一基元内部只有一组 albedo、roughness 和 normal。只要一个高斯覆盖多个像素，或者表面纹理、高光、法线在基元内部发生变化，常量材质假设就会把细节压平。SVG-IR 的改变是把每个基元改成 Spatially-varying Gaussian：在高斯局部切平面中放置多个 Gaussian vertices，每个顶点保存自己的材质和法线偏移，再在片元位置插值得到局部属性。

它的局部材质可以写成一个插值过程：
\[
  \mathbf m(\mathbf u)=\sum_k w_k(\mathbf u)\mathbf m_k,\qquad
  \mathbf m_k=(\mathbf a_k,r_k,\Delta\mathbf n_k),
\]
其中 $\mathbf u$ 是高斯局部坐标，$w_k$ 是片元到 Gaussian vertices 的插值权重。渲染流程因此接近传统图形管线。第一步是 Gaussian vertex shading：每个 vertex 根据 Disney Principled BRDF、环境光和间接光估计 radiance；第二步是 Gaussian fragment shading：像素射线命中高斯局部坐标后，在顶点属性之间插值得到该片元的 albedo、roughness 和 normal offset，再按高斯透明度混合。训练流程先用 Gaussian Surfels 初始化贴合表面的 2D 高斯，再把常量材质基元扩展为 SVG，联合优化几何、顶点材质、法线偏移和光照。

相对 GS-IR，SVG-IR 不满足于“每个高斯一组 BRDF 参数”，而是提高单个基元内部的材质分辨率；相对 TextureSplat，它不是给高斯贴显式纹理图，而是用 Gaussian vertices 和片元插值建立局部着色机制。官方仓库包含训练、重光照、PBGI、SVG splatting/rasterization 等目录，适合作为空间变化材质路线的复现实验入口。
\textbf{实现细节：} 前面所有方法都假设一个高斯内材质均匀。SVG-IR（2025.04） （Spatially-Varying Gaussian Splatting）指出这是一个根本性的表示缺陷：原始3DGS中每个高斯只能持有一组常数材质参数，在整个高斯覆盖区域内产生均匀颜色，无法捕获材质的空间变化。SVG-IR 的核心创新是空间变化高斯（SVG）表示：在每个高斯的切空间中定义 $M$ 个高斯顶点（Gaussian Vertex），类比三角网格的顶点概念，每个顶点持有独立的反照率 $a_i$、粗糙度 $r_i$ 和法线偏移 $\delta_i$。实践中 $M=4$。渲染时采用类光栅化的两阶段流程：顶点着色阶段对每个高斯顶点用基于物理的光照模型计算辐射度，得到 $M$ 个颜色 $c_i^{\{M\}}$；片段着色阶段对像素对应切空间坐标 $(u,v)$ 进行双线性插值，再经 $\alpha$-混合得到最终像素颜色：$C = \sum_{i=0}^n c_i \alpha_i T_i$，其中 $c_i = \mathrm{interp}(c_i^{\{M\}}, u, v)$。光照模型方面，SVG-IR 对每个高斯顶点求解渲染方程 $L_o(x, \omega_o) = \int_{\Omega} f(x, a, r, \omega_i, \omega_o) L_i(x, \omega_i)(\omega_i \cdot n) d\omega_i$，均匀采样 $K$ 个方向，分别计算直接光照和间接光照。间接光照通过 BVH 加速的光线追踪在已学习的辐射场中累积：$L_i^{ind}(x, \omega_i) = \sum_{l=0}^n c_l \alpha_l T_l$，可见性阈值设为 $\tau=0.8$。重光照时，SVG-IR 将训练光下的间接光照替换为一次弹射的物理间接光照，从光线命中点出发用 BRDF 重新着色，实现动态间接光照。
\textbf{方法信息：} 代码：\url{https://github.com/learner-shx/svg-ir}。
\textbf{设计动机与边界：} 相对原始 3DGS 以独立高斯和图像重建误差为主的设定，SVG-IR 把问题落在以下可实现的改动上：核心创新是SVG表示，在高斯切空间定义多个代表性位置（高斯顶点），每个顶点具有不同的材质参数（反照率、粗糙度、法线偏移）。查询高斯内任意位置的属性时，通过插值高斯顶点的属性值实现空间变化特性。SVG溅射方案分两步：高斯顶点着色计算每个顶点的光照，高斯片段着色在像素级别插值并合成。物理间接光照模型显式建模不同反弹的光传输，通过追踪光线在不同高斯间的反射计算间接光照，使间接光照在新光源下也能正确变化。使用Disney原理BRDF模型参数化材质。其设计目的不是单独增加一个模块，而是用上述表示、监督或渲染约束消除原流程在目标任务中的具体失配。提出空间变化高斯(SVG)表示，允许每个高斯内部具有空间变化的材质参数，解决了传统高斯恒定材质参数导致的渲染伪影问题。设计SVG溅射方案，类比传统图形管线的顶点/片段着色。集成基于物理的间接光照模型，实现更真实的重光照。SVG-IR在新视角合成PSNR上超越SOTA NeRF方法2.5 dB，在重光照任务上超越现有高斯方法3.5 dB，同时保持实时渲染速度。适用时还应保留论文给出的边界：SVG增加了每个高斯的参数量，存储开销增大；间接光照模型计算成本较高；对于复杂多反弹全局光照的建模仍有局限。
% METHOD: 2504.18468
\paragraph{RGS-DR（延迟反射与残差着色）（2025.04）}
RGS-DR~\cite{2504_18468} 采用 2D surfel 基元和延迟渲染处理高光物体。2D surfel 提供显式法线，比 3DGS 的最小尺度轴估计更适合高光；延迟管线则避免逐高斯先着色再混合造成的材质污染。

RGS-DR 的管线分成三遍。第一遍是 geometry/G-buffer pass：把 2D 高斯的 diffuse color、roughness、specular tint、normal、opacity 和低维 feature rasterize 到像素空间。第二遍是 lighting pass：在像素级用 split-sum IBL 近似渲染方程，环境 cubemap 的 mip level 对应不同 roughness，得到 diffuse 与 specular 分量。第三遍是 residual pass：对简化 BRDF 和单次环境光解释不了的 inter-reflection、复杂高光尾部和视角相关残差，用 G-buffer feature 加 spherical feature mip-map，再经浅层 MLP 预测残差颜色。训练时先优化几何、材质和环境光，最后冻结主要分解结果，只微调残差分支，使它提升 NVS 细节但不破坏可编辑材质。相对 Ref-Gaussian，它的重点是把“物理可解释分量”和“视觉残差分量”分开；相对 Spec-Gloss Surfels，它仍使用较常见的材质参数化，但用残差 pass 弥补高光复杂性。
\textbf{方法信息：} 代码：\url{https://github.com/gkouros/rgs-dr}。
\textbf{设计动机与边界：} 相对原始 3DGS 以独立高斯和图像重建误差为主的设定，RGS-DR 把问题落在以下可实现的改动上：RGS-DR采用2D surfel基元（显式法线表示）构建场景，每个基元具有漫反射颜色、粗糙度、镜面色调和特征向量等可学习属性。核心是延迟渲染管线：(1) G-Buffer渲染：将几何和材质属性光栅化到G-Buffer中，包括漫反射图、粗糙度图、镜面色调图等。(2) 延迟着色：基于渲染方程近似计算漫反射和镜面分量，考虑环境光照和间接光照。(3) 残差渲染通道：结合G-Buffer特征和球形特征mip-map编码物体间相互反射效果。最终渲染图像=漫反射+镜面+残差分量。方法支持高质量的重光照和材质编辑。其设计目的不是单独增加一个模块，而是用上述表示、监督或渲染约束消除原流程在目标任务中的具体失配。提出RGS-DR，采用2D surfel基元和延迟渲染管线进行高光物体的逆向渲染，引入方向性残差通道捕捉剩余视角相关效果，相比逐高斯着色方法获得更锐利的高光、更干净的材质和更好的可编辑性。适用时还应保留论文给出的边界：对极端复杂的多重反射场景建模仍有挑战；延迟渲染管线的计算开销较大；需要足够的多视角输入。
% METHOD: 2505.21238
\paragraph{3D-UIR（水下介质解耦）（2025.05）}
3D-UIR~\cite{2505_21238} 把材质/光照分解推广到水下成像。水下图像不是简单光照变化，而是后向散射、衰减和非均匀介质共同作用；直接用 3DGS 拟合会把水体颜色雾化吸收到场景外观里。3D-UIR 的改变是把清洁物体外观和水体介质效应分开建模。

它包含水下外观建模、散射介质建模和水下图像形成模型。每个高斯带外观嵌入，与视角嵌入经 MLP 产生仿射颜色变换，捕获局部颜色变化。散射模块基于投影深度估计视角特定的后向散射和衰减参数。训练时，清洁场景颜色先经过物理水下成像模型变成与观测一致的水下图像，再与真实图像计算损失。深度引导正则和自适应密度控制用于改善几何。相对 HDR/曝光方法，它处理的是传播介质而非相机响应。
\textbf{实现细节：} GI-GS（2024.10） （Global Illumination Decomposition）通过三阶段管线将全局光照显式分解为直接光照、阴影和间接光照。Baseline 方法（如 GS-IR）用球谐编码间接光照，但球谐缺乏物理约束，学到的间接光照不可重光照。GI-GS 的设计动机是：将每个物理成分建模为独立的可微操作，使分解结果可编辑、可重光照。第一阶段用标准3DGS重建几何，法线作为高斯属性优化，通过深度图导出的伪法线 $\hat{n}$ 监督渲染法线：$L_n = \|n - \hat{n}\| + \lambda_{TV} L_{TV_{\mathrm{normal}}}$，避免强制高斯退化为扁平盘的几何正则化导致的渲染质量下降。第二阶段建模直接光照：镜面分量通过 split-sum 近似分解为 $L_s = \underbrace{R}_{\mathrm{BRDF积分}} \cdot \underbrace{I_s}_{\mathrm{预滤波环境贴图}}$，漫反射分量从预卷积环境贴图 $I_d$ 查询。环境光遮蔽 $O(\mathbf{p})$ 通过从表面点向半球均匀采样射线、射线行进检测交点来计算：$O(\mathbf{p}) = 1 - \frac{1}{\pi}\int_{\Omega^{}+} V(\mathbf{p}, \omega) (\omega \cdot n) d\omega$，其中可见性 $V=1$ 当射线与表面相交。直接光照最终为 $L_{\mathrm{dir}} = (1-m)\boldsymbol{a} \cdot O(\mathbf{p}) \cdot I_d + R \cdot I_s$。第三阶段建模间接光照：从第一阶段渲染的直接光照图像 $I_{\mathrm{dir}}$ 出发，对每个表面点做路径追踪，命中点的出射辐射度直接索引 $I_{\mathrm{dir}}$ 中对应像素值：$L_i(\mathbf{p}, \omega_i) = V(\mathbf{p}, \omega_i) \cdot I_{\mathrm{dir}}(\mathbf{p}')$。最终渲染为两阶段之和 $L_o = L_{\mathrm{dir}} + (1-m)\boldsymbol{a} \cdot L_{\mathrm{ind}}$。对于复杂场景，GI-GS 进一步将屏幕空间路径追踪扩展到基于立方体图（Cubemap）的世界空间，渲染六面深度图和 RGB 图覆盖完整场景几何，使间接光照能捕获视锥外的遮挡和反射。
\textbf{设计动机与边界：} 相对原始 3DGS 以独立高斯和图像重建误差为主的设定，3D-UIR 把问题落在以下可实现的改动上：3D-UIR框架包含三个核心模块：1)水下外观建模(UAM)：为每个高斯引入外观嵌入，结合视角嵌入通过MLP进行仿射颜色变换，捕获局部变化并缓解视角依赖的颜色失真；2)散射介质建模(SMM)：基于投影深度信息估计视角特定的后向散射和衰减参数，显式建模水下光传播的物理过程；3)水下图像形成模型(UIFM)：将清洁场景颜色与介质参数通过物理方式结合，将渲染的清洁颜色转换为匹配水下拍摄图像的光度特征用于训练监督。此外，深度引导正则化优化(DRO)提供伪深度监督以改善参数估计精度，自适应密度控制动态管理高斯分布。其设计目的不是单独增加一个模块，而是用上述表示、监督或渲染约束消除原流程在目标任务中的具体失配。提出3D-UIR框架，通过物理模型将物体外观与水体介质效应解耦，解决了水下场景中非均匀散射介质对3DGS重建的干扰问题。引入外观嵌入实现后向散射和衰减的显式介质表示，提出深度引导的优化策略利用伪深度图作为监督，结合深度正则化和尺度惩罚项提升几何精度。通过水下成像模型集成外观和介质建模组件，同时实现高质量新视角合成和物理准确的场景恢复。适用时还应保留论文给出的边界：依赖水下成像模型的准确性假设；对极端水质条件（高浊度）的处理能力有限；伪深度图的质量影响几何恢复精度。
\textbf{方法信息：} DJI 参与了这篇论文；代码：\url{https://github.com/bilityniu/3D-UIR}。
% METHOD: 2506.13348
\paragraph{TextureSplat（逐基元纹理映射）（2025.06）}
TextureSplat~\cite{2506_13348} 把传统图形学的纹理映射带入 2DGS。普通 2DGS/3DGS 让每个基元只带单值法线、粗糙度或金属度，无法表达高频镜面纹理；TextureSplat 的改变是在每个 2D 高斯局部坐标系中定义材质纹理图。

2DGS 基元本身是局部平面，因此每个 splat 都有天然的 $(u,v)$ 参数化。TextureSplat 利用解析的射线-splat 交点把屏幕像素映射回该局部坐标，再从 per-primitive texture 中采样 albedo、roughness、metallic 和 tangent-space normal。切线空间法线通过基元旋转矩阵变到世界坐标后进入 deferred PBR 着色，所以一个高斯内部也能出现法线起伏和高频高光，而不必为了纹理细节继续增加基元数量。训练上先运行几何和单值材质阶段，稳定后冻结位置并打开 textured deferred PBR，优化每个基元的小纹理图；推理时把所有小纹理打包成 texture atlas，用 CUDA texture object 做硬件双线性采样。相对 SVG-IR 的 Gaussian vertices 插值，TextureSplat 更接近传统 normal mapping/texture mapping；相对 RGS-DR 的残差通道，它直接提高材质属性的空间分辨率。
\textbf{方法信息：} 代码：\url{https://github.com/maeyounes/texturesplat}。
\textbf{设计动机与边界：} 相对原始 3DGS 以独立高斯和图像重建误差为主的设定，TextureSplat 把问题落在以下可实现的改动上：该方法基于2D Gaussian Splatting（2DGS），利用其平面高斯基元自然定义的局部参数化空间进行纹理映射。核心创新在于：1）将材质属性（法线、粗糙度、金属度等）从单值属性升级为纹理图，在基元局部坐标系(u,v)中定义空间变化属性a\_k(x,y)=φ\_k(u\_k(x,y),v\_k(x,y))；2）利用2DGS的闭合射线-泼溅交点精确映射屏幕像素到局部纹理坐标，实现正确的纹理过滤；3）将切线空间法线图通过基元旋转矩阵变换到世界空间，类似传统渲染中的法线贴图；4）采用统一纹理图集（texture atlas）存储所有基元的纹理，利用GPU硬件加速纹理采样；5）基于物理的渲染模型（PBR），将分解后的材质属性与可学习环境光照结合，准确建模高反射表面的复杂光交互。该方法有效解决了现有方法因单值属性分辨率限制而无法捕捉高频镜面高光的问题。其设计目的不是单独增加一个模块，而是用上述表示、监督或渲染约束消除原流程在目标任务中的具体失配。提出将传统图形学中的纹理映射思想引入2D高斯泼溅，实现逐基元的法线图和材质纹理映射，有效解耦几何与外观表示。类似于Gouraud着色到Phong着色的演进，该方法在不增加基元数量的前提下大幅提升高频镜面反射细节的表现能力，同时通过统一纹理图集实现GPU硬件加速渲染。适用时还应保留论文给出的边界：依赖2DGS的平面基元假设，对于非平面或复杂曲面的高反射表面可能存在局限性；纹理图集的内存管理需要额外策略；方法主要针对静态场景，未涉及动态反射场景。
% METHOD: 2507.15629
\paragraph{Discretized SDF GS（离散 SDF 可重光照资产）（2025.07）}
Gaussian Splatting with Discretized SDF for Relightable Assets~\cite{2507_15629} 是本章几何-材质代表性锚点之一。它指出可重光照资产需要稳定表面，但连续 SDF 通常依赖 ray marching，高斯 splatting 则缺少 Eikonal 约束。该方法把 SDF 采样值离散地存到每个高斯里，并用 SDF 到不透明度的变换让 SDF 可以通过 splatting 渲染。

具体做法是在 2DGS/延迟着色框架中给每个基元增加一个 SDF 标量 $s_i$，并把原本直接学习的 opacity 改成由 $s_i$ 经过 SDF-to-opacity 变换得到。这样，接近零等值面的基元会产生较高不透明度，远离表面的基元贡献变小，SDF 因而可以随高斯一起 splat 到屏幕上，而不需要额外 SDF 网络或逐像素 ray marching。难点是离散 SDF 只有采样值，不能像连续场那样直接计算 $\nabla s$ 做 Eikonal 正则。为此，方法用投影一致性替代：先把高斯沿估计方向投到零等值面，再要求投影表面和 splatting 深度/法线对齐，同时用 median loss 控制 SDF-to-opacity 变换的宽度，避免所有基元都变成软雾状半透明层。

材质部分仍走高斯逆渲染常见的 deferred PBR：G-buffer 先聚合几何、法线和材质属性，再用 Principled BRDF 与 split-sum 环境光近似得到重光照结果。球面初始化用于让早期几何落在合理表面附近，减少 SDF 和 opacity 联合优化陷入局部最优。相对 RGS-DR、Ref-Gaussian 这类主要改着色和材质通道的方法，Discretized SDF GS 把“表面在哪里”也变成高斯自身属性，因此更适合作为可重光照资产的几何锚点。
\textbf{实现细节：} COREA 通过 SDF 网络和高斯的耦合来改善法线。Gaussian Splatting with Discretized SDF（2025.07） 则将 SDF 值直接编码为高斯的额外属性，避免引入额外网络表示 SDF。基于 2DGS + 延迟着色管线，核心设计是将每个高斯的不透明度从直接学习改为由 SDF 值推导：$o_i = \mathcal{T}_\beta(s_i) = 4e^{-\beta s_i}/(1+e^{-\beta s_i})^{}2$，其中 $s_i$ 为高斯的 SDF 值，$\beta$ 是全局可学习的方差控制参数——直觉上，这是一个钟形变换，SDF 值为零（恰好在表面上）时不透明度最高，SDF 绝对值越大（离表面越远）不透明度越低。为加速联合优化 $\beta$ 和 $s_i$ 的收敛，引入中值损失：监控所有高斯 $|s|_{med}$（SDF 绝对值的中位数），要求变换后的中值不透明度 $o_{med} \leq 0.5$，推导出引导值 $\beta_{med} = -(3-2\sqrt{2})|s|_{med}$，中值损失 $L_\beta = \max(\beta_{med} - \beta, 0)$。由于离散 SDF 无法直接施加 Eikonal 条件 $|\nabla f|=1$（仅法线方向可用），提出投影一致性损失：将高斯投影到 SDF 零水平集 $\mu_{proj} = \mu_i - s_i \frac{\nabla f_i}{|\nabla f_i|}$，约束投影点深度与 $\alpha$-混合 深度之差 $\Delta_{proj}^i = |D_{agg} - D_{proj}|$，当 $\Delta_{proj}^i \leq \tau$ 时施加 $L_{proj} = \frac{1}{N}\sum_{i}^{N}\Delta_{proj}^i$（超阈值的高斯为遮挡面或不可微区域，直接丢弃）。此外，球面初始化（单位球上均匀采样点）避免初始几何陷入局部最优。输出为可重光照的资产（mesh + 材质贴图）。
\textbf{设计动机与边界：} 相对原始 3DGS 以独立高斯和图像重建误差为主的设定，Gaussian Splatting with Discretized SDF for Relightable Assets 把问题落在以下可实现的改动上：核心方法：(1)离散化SDF：在每个高斯上编码一个SDF采样值作为额外属性（除PBR属性、缩放和四元数外），通过SDF到不透明度变换o\_i=T\_σ(s\_i)=4e\^{}(-σs\_i)/(1+e\^{}(-σs\_i))\^{}2从SDF值推导不透明度，而非直接学习，实现通过泼溅渲染SDF并避免光线行进；(2)投影一致性损失：由于离散表示无法应用基于梯度的Eikonal损失等约束，提出将高斯投影到SDF零等值面上，并强制与泼溅得到的表面对齐，证明是Eikonal损失的近似；(3)球面初始化：受NeuS/VolSDF启发，提出针对高斯泼撒的球面初始化策略，避免几何陷入局部最优。基于2DGS和延迟着色构建框架，使用Principled BRDF和split-sum近似作为着色模型。其设计目的不是单独增加一个模块，而是用上述表示、监督或渲染约束消除原流程在目标任务中的具体失配。提出离散化SDF，将连续SDF以离散方式编码在每个高斯内，通过SDF-不透明度变换实现通过泼溅渲染SDF，避免光线行进计算。提出投影一致性损失替代不可行的梯度约束，无需额外内存或复杂手动优化，在重光照质量上优于现有高斯基逆向渲染方法。适用时还应保留论文给出的边界：离散SDF分辨率受限于高斯密度；投影一致性损失为Eikonal损失的近似，可能不如精确约束；球面初始化假设可能不适用于所有场景。
\textbf{方法信息：} 代码：\url{https://github.com/nk-cs-zzl/discretizedsdf}；GitHub 关注度较高，约 109 stars。



% METHOD: 2510.02069
\paragraph{Spec-Gloss Surfels（高光物体的参数化选择）（2025.10）}
Spec-Gloss Surfels~\cite{2510_02069} 的增量在 BRDF 参数化。常见 Metallic-Roughness 参数化会把漫反射和镜面反射耦合，逆渲染时容易让金属度和底色互相解释；Spec-Gloss 参数化直接用独立镜面反射色和光泽度，减少高光物体歧义。

方法仍以 2DGS surfel 和 deferred shading 为基础，但把材质变量改成 diffuse albedo、独立的 specular reflectance $F_0$ 和 glossiness/roughness，而不是用 metallic 把 $F_0$ 绑到 base color 上。像素级 G-buffer 先得到 albedo、roughness、$F_0$、normal 和间接光相关项，再用可微的 prefiltered cubemap 做 split-sum IBL。早期优化阶段引入 StableNormal 和 StableDelight 一类扩散先验：法线先验帮助表面朝向收敛，diffuse prior 帮助把高光从反照率中剥离。环境光采用 coarse-to-fine cubemap，从低分辨率逐步上采样到高分辨率；negative-only clipping 只抑制不合理负值，不压低 HDR 峰值，因此镜面高光不会被训练过程“磨平”。相对 RGS-DR，它更强调材质参数化和先验约束；相对 GlossyGS，它的锚点价值在于把可编辑材质、HDR 环境光和扩散先验整理成一条可复现的 2DGS 逆渲染流程。
\textbf{实现细节：} Spec-Gloss Surfels（2025.10） 基于 2DGS 面元表示，引入了 Specular-Glossiness（简称 SG）材质参数化来进行逆渲染。当前主流方法大多采用 Disney BRDF 的 Metallic-Roughness 参数化方式，其镜面反射率 $F_0$ 由公式 $(1-m)\cdot 0.04 + m\cdot b$ 隐式决定。从直觉上说，这种参数化把"这个表面是金属还是电介质"和"它的基础颜色是什么"混在一个参数 $b$ 中，在逆渲染过程中优化器很难区分二者。在渲染方面，Spec-Gloss Surfels 采用 split-sum 近似来计算出射辐射度：$L_o = \rho_d L_d(n) + (F_0 \beta_1 + \beta_2) L_s(v,n,r)$，其中 $L_d$ 为漫反射辐照、$L_s$ 为预过滤镜面辐照、$\beta_1,\beta_2$ 为预计算 2D LUT 返回的 Fresnel 系数。环境光用可学习 HDR cubemap 表示，训练时从低分辨率 $Res^{}1$ 渐进上采样至 $Res^K$（每阶段分辨率翻倍），避免高分辨率光照过早吸收材质误差。法线先验来自 StableNormal，漫反射颜色先验来自 StableDelight 扩散模型（在训练初期 $T_{dp}$ 步后关闭），两者共同抑制几何与材质歧义。
\textbf{方法信息：} 代码：\url{https://github.com/gkouros/specgloss-gs}。

一个需要持续关注的技术候选是 3DSS~\cite{2605_05876}。它不是在 3DGS/2DGS 上再加材质通道，而是回到 EWA Surface Splatting 本身：用有向 surfel 的重构核定义连续表面，并把 surface separation 写成按深度排序 primitive stream 的 interval merging；每个 surface layer 的 opacity 来自累积 EWA 权重，因此轮廓抗锯齿和可见性梯度都从同一覆盖度模型中产生。随后，方法先在 surfel 上做 forward microfacet shading，再进行 object-space signal reconstruction 和多层 surface compositing，避免体渲染式混合把不同表面的法线、材质和辐射混在同一个像素里。密度控制也随之改变：分裂和剪枝根据 surfel 是否足以覆盖局部表面来做，最后的有向点集可接 Screened Poisson 等网格重建流程。当前 arXiv 页面没有确认完整官方代码；此前检索到的 \url{https://github.com/fatma18F/3DSSL-SS24_Generalizable3DHeadReconstruction} 是 HeadSplat/课程项目，不属于 3DSS。因此这里不把 3DSS 作为已验证可信方法展开，只把它列为可微表面 splatting 逆渲染路线的后续重点关注对象。
\textbf{设计动机与边界：} 相对原始 3DGS 以独立高斯和图像重建误差为主的设定，Spec-Gloss Surfels and Normal-Diffuse Priors for Relightable Glossy Objects 把问题落在以下可实现的改动上：方法基于2D Gaussian Splatting构建场景表示，采用延迟着色(deferred shading)管线：(1)Specular-Glossiness参数化：选择Spec-Gloss参数化替代主流的Metallic-Roughness参数化，分析发现MR参数化中F0=(1-m)*0.04+m*b将漫反射和镜面反射耦合，导致逆渲染中的歧义。Spec-Gloss通过独立的镜面反射色F0和光泽度参数解耦漫反射和镜面分量，更适合逆渲染。(2)扩散先验：利用基于扩散模型的法线和漫反射色先验指导早期优化，减少材料-照明的歧义。(3)粗到细环境光优化：渐进式上采样环境cubemap加速收敛，并通过negative-only裁剪保留高动态范围镜面反射。整体基于2DGS的显式法线和视图一致性，更适合高光效果的建模。其设计目的不是单独增加一个模块，而是用上述表示、监督或渲染约束消除原流程在目标任务中的具体失配。提出将Specular-Glossiness BRDF参数化集成到2DGS延迟着色框架中，配合法线和漫反射色的扩散先验以及粗到细环境光优化，实现高质量的光泽物体重照明。适用时还应保留论文给出的边界：环境光优化仍需多视角输入；极端高光（如镜面反射）的处理仍有挑战；依赖2DGS的表面质量
% METHOD: 2510.11387
\paragraph{MaterialRefGS（多视图一致材质推断）（2025.10）}
MaterialRefGS~\cite{2510_11387} 从多视图一致性重新审视反射建模。高反射物体中，同一个材质点在不同视角下颜色变化大，若不显式约束材质图一致，优化会把光照和材质混在一起。方法基于 2DGS 和延迟 PBR 管线，要求高斯在多视图中产生一致材质图。

它以 2DGS surfel 为基础，先把反照率、金属度、粗糙度、法线等属性渲染成像素级 G-buffer，再用简化 Disney BRDF 和 split-sum 近似完成 deferred PBR。关键不是换一个 BRDF，而是给材质推断增加多视图约束：同一组 2D 高斯从不同相机观察时，输出的材质图应当一致；跨视图光度变化大的区域会被标记为高反射候选，作为 specular strength 的先验；物体间遮挡产生的间接光照则通过 2DGS ray tracing 建模，而不是简单塞进残差颜色。相对 Ref-Gaussian，它更强调“材质图要跨视图稳定”；相对 RGS-DR，它更少依赖事后残差，更直接把反射强度和间接辐射写进优化目标。
\textbf{实现细节：} Reflections Unlock 通过分支解耦实现了反射与透射的分离，但基于 PBR 框架的方法还面临一个更深层的问题：同一表面点从不同视角渲染出的材质参数互相矛盾。MaterialRefGS（2025.10） 在原始 2DGS 的基础上修改了材质参数的约束机制——用跨视角材质一致性约束替代无约束的自由优化——Splatting 流程和 BRDF 着色模型保持不变。动机。PBR 框架假设视角无关的材质属性（反照率、粗糙度、金属度）可以通过 BRDF 解耦出视角相关的镜面反射，但在实际优化中，不同视角下高斯覆盖的像素不同、$\alpha$-混合权重各异，同一表面点从不同视角渲染出的材质图往往严重不一致——BRDF 难以从这些矛盾的材质观测中推断出全局物理反射效果。MaterialRefGS 以 2DGS 为基线表示，为每个高斯附加四个材质属性：漫反射颜色 $c_d^{}3$、反照率 $a$、金属度 $m$、粗糙度 $r$。通过 $\alpha$-混合 分别渲染出对应的材质图 $\hat{C}_d, \hat{A}, \hat{M}, \hat{R}$。
\textbf{方法信息：} 引用很多，Semantic Scholar 约 20 次。



% METHOD: 2603.27516
\paragraph{SGS-Intrinsic（稀疏视角室内逆渲染）（2026.03）}
SGS-Intrinsic~\cite{2603_27516} 面向稀疏视角室内场景。逆渲染在稀疏视角下最容易失败的是几何和语义基础不稳：阴影会被误认为材质，语义不一致区域会产生破碎反照率。SGS-Intrinsic 的改变是先构建几何一致的语义高斯场，再进行材质-光照分离。

具体流程分两步。第一步先用 VGGT 从稀疏输入构建较密集的场景点云，再把 StableNormal 的法线预测和 LSEG 语义特征作为监督写入高斯场；训练视图和虚拟视图之间还做语义一致性约束，SAM 区域掩码用于把同一物体区域内的语义特征拉到一致，避免高斯只在少数训练视角上过拟合纹理。第二步冻结或稳定几何基础后，为高斯增加法线、反照率、粗糙度等逆渲染属性，并用远场环境光加近场 Spherical Gaussian 点光源描述室内照明。光照不变材质约束要求同一语义/材质区域在不同视角和光照解释下保持相近反照率，去阴影模块则显式减少 cast shadow 被写入 albedo 的风险。相对 DeferredGS 的通用 G-buffer 管线，SGS-Intrinsic 的关键不在延迟着色本身，而在把稀疏视角下最缺的几何、法线和语义先验先补稳。
\textbf{实现细节：} SGS-Intrinsic（2026.03） （Semantic-Invariant）与 GAINS 解决类似的稀疏视角问题，但更侧重于利用语义先验改善室内场景的材质分解。Baseline 方法在稀疏视角下用 RGB 重建损失训练高斯，容易过拟合到已观测视角，几何不可靠，且光照与材质的解耦陷入局部最优——学到的反照率常将光照效果"烘焙"进去。SGS-Intrinsic 的核心观察是：室内场景中同一语义类别（如"桌面""墙面"）的物体通常具有相似材质属性，因此可通过语义一致性约束大幅减少独立参数量并正则化解耦。方法分两阶段。阶段 I 使用预训练模型提供额外监督：法线引导损失 $L_{\mathrm{normal}} = 1 - \hat{n}^T n_m$（来自 StableNormal，$\lambda_1=0.05$）；语义监督利用 LSEG 编码器提取像素级 512 维特征 $F$，Qwen3VL 识别物体类别后经 CLIP 生成文本嵌入 $T \in \mathbb{R}^{M \times 512}$，余弦相似度构建语义标签图 $S_{i,j} = \arg\max_m \cos(F_{i,j}, T_m)$，语义损失为 $L_{\mathrm{semantic}} = L_{\mathrm{cos}}(\hat{F}, F) + L_{\mathrm{ce}}(\hat{S}, S)$（$\lambda_2=0.3$）；训练-伪视角语义一致性约束 $L_{\mathrm{corresponding}}$ 在 SAM 生成的掩码区域内执行多数投票，强制语义一致。阶段 II 在固定几何上进行逆渲染：混合光照模型将入射光分解为高频球面高斯点光源 $L_i^{\mathrm{SGM}}$（$N_{\mathrm{sg}}$ 个 SG 叶瓣）和低频可学习环境 cubemap $L_i^{\mathrm{env}}$（$6 \times 512 \times 512 \times 3$）；轻量级去阴影模块用多级哈希编码提取空间特征，结合点光源的相对方向和距离，通过 MLP 预测可见性 $V_{ji}$ 以显式建模投射阴影。关键的材质一致性正则化包含两个损失：重光照一致性 $L_{\mathrm{relight}}$（$\lambda_4=0.4$）在固定视角施加 $N$ 次光照扰动，用 MSE 和多尺度梯度损失强制反照率不变；跨视角语义-材质一致性 $L_{\mathrm{cross_semantic}}$（$\lambda_3=0.4$）假设同类语义区域材质相似——$L_{\mathrm{cross_semantic}} = \sum_{p=1}^P (\|\bar{F}_p^{*} - \bar{F}\|_2 + \|\bar{A}_p^{*} - \bar{A}\|_2)$，其中 $\bar{F} = \sum(F \odot M) / \sum M$ 为掩码内平均语义特征，$\bar{A}$ 为平均反照率。整个训练在单张 RTX 4090 上约 40 分钟（阶段 I 7000 次 + 阶段 II 3000 次迭代）。
\textbf{设计动机与边界：} 相对原始 3DGS 以独立高斯和图像重建误差为主的设定，SGS-Intrinsic 把问题落在以下可实现的改动上：SGS-Intrinsic采用两阶段框架：第一阶段（几何与语义重建）：使用VGGT构建密集场景点云，融合StableNormal法线先验和LSEG语义特征作为额外监督信号。引入训练视图与伪视图间的语义一致性约束，通过SAM生成的掩码进行区域级语义均匀化，确保语义场的几何一致性。第二阶段（逆向渲染）：结合近场点光源和远场环境光的混合光照模型，使用材质先验（如阿拉bedo的平滑性）进行材质-光照分离。引入光照不变材质约束减少投射阴影影响，并设计去阴影模型增强材质恢复鲁棒性。每个3D高斯额外存储法线、反照率、粗糙度等逆向渲染属性。其设计目的不是单独增加一个模块，而是用上述表示、监督或渲染约束消除原流程在目标任务中的具体失配。提出首个面向稀疏视角室内场景的3DGS逆向渲染框架。通过构建几何一致的语义高斯场为逆向渲染提供可靠基础，结合混合光照模型和材质先验实现材质-光照分离，引入光照不变材质约束和去阴影模型减轻阴影对材质恢复的影响，在几何重建和逆向渲染质量上均优于现有3DGS方法。适用时还应保留论文给出的边界：依赖预训练模型的先验质量，混合光照模型假设可能不适用于所有室内场景，计算开销较高。
\textbf{方法信息：} Amazon 参与了这篇论文；代码：\url{https://github.com/grumpysloths/sgs_intrinsic.github.io}。



% METHOD: 2604.10994
\paragraph{LumiMotion（利用动态线索的逆渲染）（2026.04）}
LumiMotion~\cite{2604_10994} 的局部锚点是 2DGS 逆渲染，但它利用了静态场景方法忽略的线索：运动物体会让同一表面在不同光照和遮挡状态下被多次观察，从而帮助区分材质和光照。普通单状态逆渲染容易把阴影或高光吸收到反照率里；LumiMotion 用场景动态提供额外约束。

具体流程分两阶段。第一阶段仍是动态 2DGS 重建：用一组 canonical 2D Gaussian surfel 表示场景表面，再用时间条件 MLP 预测位置、旋转和颜色变化。这里故意不让变形网络随意改变 opacity 或 scale，因为方法希望动态物体通过空间运动暴露同一材质在不同光照/遮挡下的观测，而不是让网络把阴影解释成“某些高斯出现或消失”。论文还加入静态-动态模糊分离约束，使真正运动的区域可以变形，背景和静态表面则尽量保持稳定；这一步的输出是带稳定法线的 2DGS 几何和一个已训练好的动态变形网络。

第二阶段冻结几何和变形网络，改做逆渲染。此时不再直接渲染高斯颜色，而是为每个 surfel 优化 albedo 和 roughness，并用简化 Disney BRDF、环境光和 ray tracing 估计入射光，最后由渲染方程生成观测颜色。动态区域在训练中相当于提供了同一表面被阴影扫过、被物体遮挡或在不同直接光照下出现的多次观测，因此 albedo 更不容易把 cast shadow 烘焙进去。相对 SGS-Intrinsic 的语义/法线先验，LumiMotion 的额外先验来自运动本身；相对 GI-GS/RGS-DR 这类静态逆渲染，它把“动态场景很难重光照”反过来变成材质-光照分离的监督来源。官方仓库包含分阶段训练、数据准备和 surfel tracer/rasterizer 依赖，说明该方法不仅是概念验证，而是一套可运行的动态逆渲染流水线。
\textbf{实现细节：} 前面所有方法都假设场景是静态的。LumiMotion（2026.04） （Improving Gaussian Relighting with Scene Dynamics）解决静态光照假设下动态场景的重光照问题：先前方法假设场景中光照恒定，但真实动态场景中光照和物体同时变化（如人在室内走动时的阴影变化）。Baseline 2DGS 提供准确的表面法线（对 BRDF 计算至关重要），但无法建模时间维度的变化。LumiMotion 的两阶段管线如下。阶段 I 学习动态几何：基于 2DGS，用 MLP 预测每个时间步 $t \in [0,1]$ 下高斯的平移 $\Delta \mu \in \mathbb{R}^{}3$、旋转 $\Delta r \in \mathbb{R}^{}3$ 和颜色变化 $\Delta c \in \mathbb{R}^{}3$——$\Delta \mu, \Delta r, \Delta c = \mathrm{MLP}(\gamma(t), \gamma(\mu))$，其中 $\gamma$ 为位置编码。关键设计是静态-动态模糊分离：为每个高斯引入辅助变量 $P$，用 Binary Concrete 分布采样（温度 $T=0.5$）得到连续松弛变量 $\tilde{P} \in [0,1]$，仅动态高斯（$\tilde{P}$ 接近 1）才施加平移和旋转变换——$\mu' = \mu + \tilde{P} \cdot \Delta \mu$，$r' = r + \Delta r$。动机是：移动阴影在无纹理区域可被错误地解释为高斯的颜色变化或消失，静态-动态分离确保动态变化仅作用于真正运动的物体，这对阶段 II 的反照率估计至关重要。颜色变化采用乘性模型 $c' = c(1 - \Delta c)$ 模拟光影效果。阶段 II 在冻结的几何和变形网络上进行逆渲染：每个 2D 高斯分配漫反射反照率 $A$（用阶段 I 的规范颜色 $c$ 初始化）和粗糙度 $\rho$，两者对所有时间步保持恒定（颜色变化完全由渲染方程解释，$\Delta c$ 不再使用）。环境光照 $L_{\mathrm{env}}$ 用可优化图像表示。渲染方程在光栅化后逐像素求解（而非逐高斯），采用简化 Disney BRDF（反照率 $A$ + 粗糙度 $\rho$），入射辐射度 $L_i = V(\omega_i, x) L_{\mathrm{env}}(\omega_i) + L_{\mathrm{ind}}(\omega_i, x)$ 通过 2D 高斯光线追踪计算可见性 $V$ 和间接光照 $L_{\mathrm{ind}}$。
\textbf{方法信息：} Microsoft 参与了这篇论文；代码：\url{https://github.com/joaxkal/lumimotion}。
\section{反射与镜面建模}

反射路线的核心问题是，镜子、玻璃和高光会把别处的场景带到当前表面。原始 3DGS 往往把反射当成虚假几何、漂浮点或高阶视角颜色，导致新视角不一致。方法大致分三类：虚拟相机解释镜面，延迟反射解释高光，增强核或潜特征补偿 SH 无法表达的尖锐视角相关外观。


% METHOD: 2404.01168
\paragraph{Mirror-3DGS（平面镜虚拟视点）（2024.04）}
Mirror-3DGS~\cite{2404_01168} 直接利用平面镜成像原理：镜中虚像等价于从镜面后方的对称视点观察真实世界。原始 3DGS 会把镜中内容误当成镜后虚假几何；Mirror-3DGS 的改变是在渲染时显式构造镜像相机。

训练时，方法为每个高斯增加 mirror attribute，并用镜面 mask 监督这一属性；第一阶段用固定颜色替换镜面内容，使高斯先学习真实场景和镜面位置，避免把虚像也长成一组几何。第二阶段筛选 mirror attribute 高的高斯拟合平面方程，并加入平面一致性和深度约束稳定平面估计。渲染时分两路：非镜面区域由原相机正常 3DGS 渲染；镜面区域先根据平面反射构造虚拟相机，从镜后对称视点渲染真实场景，再用镜面 mask 合成。相对 Gaussian Splatting in Mirrors，它的核心更直接地依赖镜面属性和平面镜一次反射；后者则把精力放在虚拟相机参数的进一步优化。
\textbf{实现细节：} Mirror-3DGS（2024.04） 修改了原始 3DGS 的渲染管线——用虚拟镜像视角的双路渲染替代单视角 $\alpha$-混合——高斯的几何参数和 SH 颜色保持不变。动机。镜子的特殊性在于反射内容构成一幅完整的"虚拟图像"，Baseline 3DGS 的 $\alpha$-混合沿光线累积颜色，容易在镜子后方生成错误的高斯——这些高斯违反物理规律（物体不可能在镜子内部），严重损害渲染质量。核心思路来自几何光学的镜像成像原理：镜面呈现的虚拟图像等价于从镜子后方的对称位置观察真实世界，因此无需昂贵的光线追踪即可实现物理正确的镜面反射。Mirror-3DGS 的实现分为三个关键步骤。首先，为每个高斯添加可学习的镜面属性 $m \in [0,1]$，表示该高斯属于镜子的概率，通过与 $\alpha$-混合 相同的方式渲染出镜面掩码： 使用数据集提供的真值镜面掩码 $M_{gt}$ 通过 $\mathcal{L}_{mask} = \mathcal{L}_1(M, M_{gt})$ 监督学习。然后，基于高镜面概率的高斯点，用改进的 RANSAC 算法剔除噪声后拟合镜面平面方程 $\Pi = (\mathbf{n}_\Pi^{}\top, d)$，其中 $\mathbf{n}_\Pi$ 为法向量，$d$ 为平面到原点的距离。平面约束损失为： 确保所有镜面高斯都落在同一平面上。最后，基于平面方程构造虚拟镜像视角：将原始相机位姿 $P_o \in \mathbb{R}^{4 \times 4}$ 通过镜像变换矩阵 $T_m$ 映射到镜后位置 $P_m = T_m P_o$。掩码融合与两阶段训练。渲染时从原始视角得到图像 $C_{P_o}$（保留镜外区域），从虚拟视角得到图像 $C_{P_m}$（保留镜内区域），最后通过掩码融合： 其中 $\odot$ 为逐元素乘法。训练采用两阶段策略。第一阶段用红色填充真值图像中的镜面区域（消除镜面内容干扰），联合优化高斯表示和镜面平面方程，损失为 $\mathcal{L}_{stage}^{}1 = \lambda_{mask}\mathcal{L}_{mask} + \mathcal{L}_{rgb}^{}1 + \lambda_{depth}\mathcal{L}_{depth} + \mathcal{L}_{plane}$。第二阶段固定镜面平面方程，使用完整的未遮挡真值图像 $C_{gt}$ 监督融合结果，损失为 $\mathcal{L}_{stage}^{}2 = \lambda_{mask}\mathcal{L}_{mask} + \mathcal{L}_{rgb}^{}2$。分阶段训练的动机是：先让高斯学会场景的几何结构和镜面位置，再通过虚拟视角补充镜面反射的细节。
\textbf{方法信息：} 第三方完整实现：\url{https://github.com/gragonfly/mirror-3dgs-my}。
\textbf{设计动机与边界：} 相对原始 3DGS 以独立高斯和图像重建误差为主的设定，Mirror-3DGS 把问题落在以下可实现的改动上：方法核心创新是利用「移动到镜面后方观察世界」的镜像成像原理：镜面上呈现的虚像等价于从镜面后方观察真实世界。具体步骤包括：1）为3D Gaussians引入镜面属性标识，估计镜面平面方程；2）从原始视点渲染图像，保留镜面外区域；3）根据镜面平面计算镜像视点，从镜像视点渲染获得镜面反射内容；4）通过镜面掩码将两幅图像合成最终结果。方法避免了Mirror-NeRF所需的耗时光线追踪，直接利用3DGS的高效光栅化。使用分割工具自动检测镜面区域，估计镜面平面参数。支持多镜面场景和不同镜面形状。其设计目的不是单独增加一个模块，而是用上述表示、监督或渲染约束消除原流程在目标任务中的具体失配。提出Mirror-3DGS框架，通过在3DGS中引入镜面属性并利用平面镜成像原理，从镜面后方模拟镜像视点来增强场景渲染的真实感。无需耗时的光线追踪即可实现物理准确的镜面反射效果。在合成和真实场景上实时渲染新视角，在镜面区域超越Mirror-NeRF。适用时还应保留论文给出的边界：需要镜面区域的检测和平面估计，对非平面镜（如凸面镜/凹面镜）不适用；镜面检测依赖分割工具可能不够鲁棒；多镜面互相反射场景处理复杂；假设镜面为理想反射。
% METHOD: 2404.18454
\paragraph{3DGS with Deferred Reflection（两遍延迟反射）（2024.04）}
3D Gaussian Splatting with Deferred Reflection~\cite{2404_18454} 是反射路线的代表性锚点。原始 3DGS 用 per-Gaussian spherical harmonics 表达视角相关颜色，适合中低频外观，但镜面反射是高频环境查询；如果仍让 SH 直接拟合，模型往往会生成虚假高斯来“画出”反射图案，几何和外观都被污染。该方法的改变是把镜面反射从单遍颜色 splatting 中拆出来，改成两遍 deferred shading。

第一遍仍执行 3DGS splatting，但除基础颜色外，还把每个高斯的法线和反射强度 splat 到屏幕空间。法线来自高斯椭球最短轴，并在渲染时翻转到朝向相机；反射强度是每个高斯学习的标量。第二遍在像素级做反射：用屏幕空间法线和视线方向得到反射方向，查询可学习 environment map，然后按反射强度把基础颜色和反射颜色混合。这样，昂贵的排序和 splatting 权重仍沿用 3DGS，新增的反射计算只发生在屏幕像素上，运行帧率接近 vanilla 3DGS。

训练难点是法线。环境贴图查询的梯度来自双线性纹理过滤，只有反射方向已经接近正确 texel 时才有稳定信号；错误法线很难从颜色损失中自己转正。Deferred Reflection 因此引入 normal propagation：先在短暂的 view-independent warm-up 后打开反射强度和环境图优化，找出反射强度较高、法线已近似正确的高斯；训练中周期性提高这些高斯的 opacity/反射强度并扩展其长轴，使它们与邻近高斯共享像素，邻居就能通过同一像素的 deferred normal 获得有意义梯度。相对 forward shading 中每个高斯独立查询反射，deferred shading 让多个高斯在同一像素的 blended normal 上相互传递法线信息，这正是它能稳定处理镜面/高光反射的核心。
\textbf{实现细节：} 3D Gaussian Splatting with Deferred Reflection（2024.04） 修改了原始 3DGS 的渲染管线——用两阶段延迟着色（deferred shading）替代单遍 SH 颜色 Splatting——高斯的位置、协方差、不透明度等几何参数保持不变。动机。标准 3DGS 用球谐（SH）编码视角相关颜色，将几何信息和外观信息混在同一组参数中——当场景包含镜面反射表面时，SH 的表达能力不足以同时编码漫反射和反射内容，导致重建质量受限。直觉上，延迟着色的思路是"先收集几何信息，再统一着色"——就像画家先画好线稿再上色，而不是边画边决定最终颜色。第一阶段是标准高斯 Splatting，但除了颜色 $c_i(v)$ 外，同时将每个高斯的法线向量 $\mathbf{n}_i$（取椭球最短轴）和镜面反射强度 $r_i$（标量）作为额外属性一起 Splat 到屏幕空间： 其中 $G(\boldsymbol{\mu}_i, v)$ 是标准 Splatting 的 $\alpha$-混合 权重，复用了已经计算好的排序结果，因此额外开销极小。法线 $\mathbf{n}_i$ 根据视角方向翻转以确保始终朝向相机。第二阶段：延迟反射着色。在屏幕空间中组合最终像素颜色： 直觉上，$R(v)$ 是逐像素的反射强度图：$R(v) \approx 0$ 处保留原始漫反射颜色 $C(v)$；$R(v) \approx 1$ 处完全由环境贴图（Environment Map）$E$ 提供反射颜色。$E$ 在反射方向上用双线性插值查询。训练策略。关键设计在于：环境贴图 $E$ 完全从最终渲染损失端到端训练，与高斯参数解耦。反射强度 $r_i$ 初始化为 $0$，训练初期先进行几千次迭代的视角无关优化（SH 限制在 0 阶，即常数颜色）。
\textbf{锚点细化：} 3DGS with Deferred Reflection：针对球谐将几何与外观混在同一组参数、表达镜面反射时能力不足的局限，用两阶段延迟着色（deferred shading）替代单遍 SH 颜色 Splatting。延迟着色的直觉类似画家"先打线稿再上色"：第一阶段只把每个高斯的法线、镜面强度等几何属性 Splat 到屏幕空间，第二阶段才在像素级统一查询环境贴图得到反射颜色——这样几何参数完全不动，手术只发生在外观环节，避免了几何与材质互相干扰。
\textbf{方法信息：} 代码：\url{https://github.com/gapszju/3dgs-dr}；GitHub 关注度较高，约 187 stars；引用很多，Semantic Scholar 约 81 次。



% METHOD: 2409.05868
\paragraph{Latent-SpecGS（潜特征镜面外观）（2024.09）}
SpecGaussian with Latent Features~\cite{2409_05868} 认为低阶 SH 不足以表达尖锐镜面和各向异性外观。它为每个高斯嵌入神经描述符，并通过解码器生成漫反射、镜面和视角掩码。

渲染时，Feature Gaussian Splatting 先输出漫反射特征、视角依赖特征和掩码特征图。UNet 风格网络解码漫反射颜色，紧凑 CNN 加视角编码解码镜面颜色；视角掩码决定镜面分量如何叠加到漫反射上。法线也由潜描述符估计。相对 Augmented Radiance Field 的增强核，它把复杂外观交给 2D 特征图和神经解码器，而不是增加视角依赖 opacity。
\textbf{方法信息：} 代码：\url{https://github.com/marcwangzhiru/speclatentgs}。
\textbf{设计动机与边界：} 相对原始 3DGS 以独立高斯和图像重建误差为主的设定，SpecGaussian with Latent Features 把问题落在以下可实现的改动上：方法首先通过COLMAP生成稀疏点云初始化3D Gaussians，每个Gaussian赋予神经描述符（包含漫反射特征、视角依赖特征、视角掩码特征和法向量）。通过Feature Gaussian Splatting生成特征图（漫反射、视角依赖、掩码）。使用UNet风格网络解码漫反射颜色，紧凑CNN加视角编码解码镜面颜色。借鉴Cook-Torrance模型，将镜面颜色乘以视角掩码后叠加到漫反射颜色上。法向量通过潜在描述符计算得到。其设计目的不是单独增加一个模块，而是用上述表示、监督或渲染约束消除原流程在目标任务中的具体失配。提出Latent-SpecGS，在每个3D Gaussian中嵌入通用潜在神经描述符以增强外观和几何表示。设计双分支CNN解码器分别生成漫反射和镜面反射颜色图，通过视角相关掩码自适应融合。受Cook-Torrance模型启发，有效处理复杂镜面反射场景，扩展了3D-GS处理复杂光照的能力。适用时还应保留论文给出的边界：神经描述符增加了每个Gaussian的存储开销；双分支解码器增加了训练和推理的计算量。
% METHOD: 2410.01614
\paragraph{Gaussian Splatting in Mirrors（虚拟相机优化）（2024.10）}
Gaussian Splatting in Mirrors~\cite{2410_01614} 同样采用镜像视点思路，但强调虚拟相机优化。镜面平面估计稍有误差，反射内容就会在新视角中漂移；直接用 RANSAC 平面可能不够。

方法先从 3DGS 渲染得到深度、法线和镜面区域 mask，选择镜面像素足够多的视图，把这些像素按深度和相机内外参回投到世界坐标，再用区域平均法线和 RANSAC 拟合镜面平面。平面确定后，真实相机可按镜面对称变换得到 virtual camera，用它解释镜中反射内容。由于深度、法线和 mask 都可能有误差，论文继续把虚拟相机位姿作为可优化参数，用反射区域的 RGB 损失微调，使镜中内容在多视角下对齐。相对 Mirror-3DGS 的 mirror attribute + 平面拟合路线，它更像在几何估计后增加一层相机校准；局限是新视角镜面 mask 仍需要额外来源，仓库 README 也明确提示这一点。
\textbf{方法信息：} 代码：\url{https://github.com/rzhevcherkasy/bmvc24-gsim}。
\textbf{设计动机与边界：} 相对原始 3DGS 以独立高斯和图像重建误差为主的设定，Gaussian Splatting in Mirrors 把问题落在以下可实现的改动上：方法分三步：(1)镜面平面估计：从3DGS获取深度图、法线图和像素级镜面区域掩码，随机选择包含大镜面区域的视角，计算世界坐标和法线值，通过RANSAC拟合最佳平面方程；(2)虚拟相机构建：根据估计的镜面平面方程，在对称位置放置虚拟相机来解释镜面反射内容；(3)虚拟相机优化：针对镜面平面估计不完美的问题，提出简单有效的虚拟相机优化方法增强反射质量。收集了包含三个真实世界场景的新镜面数据集用于评估。其设计目的不是单独增加一个模块，而是用上述表示、监督或渲染约束消除原流程在目标任务中的具体失配。提出将镜面反射建模为基于物理的虚拟相机的新方法，通过优化虚拟相机实现高质量多视角一致的反射渲染，训练时间相比先前SOTA大幅减少。适用时还应保留论文给出的边界：依赖镜面区域的准确检测；对非平面反射面（如曲面镜）处理有限；镜面平面估计精度影响最终反射质量
% METHOD: 2410.02764
\paragraph{Flash-Splat（闪光线索的 3D 反射分离）（2024.10）}
Flash-Splat~\cite{2410_02764} 处理玻璃等透明反射场景中的传输层和反射层分离。传统 flash/no-flash 反射分离要求同一相机位置拍一张闪光、一张无闪光；只要相机轻微移动，二者相减就会产生严重错位。Flash-Splat 的改变是把“配对差分”提升到 3D：用户可以先采集一组多视角闪光图像，再采集一组多视角无闪光图像，两组视角只需覆盖相近区域，不需要逐帧对齐。

系统的输入是两组未配对多视角图像和 COLMAP 相机，内部显式维护 transmission-with-flash、transmission-without-flash、reflection 和 reflection fraction 等 3D 表示，并利用 3DGS 的新视角合成能力构造两类 pseudo-pair：一类是在某个已拍闪光视角合成对应无闪光图像，另一类是在 3D 空间分别重建闪光/无闪光表示后做差。闪光主要增强玻璃后的传输层，对反射层影响较小，这个物理差异被写成 flash prior 损失来约束两层分离；训练时还加入普通 RGB 重建损失，使分离后的两层重新合成后仍能渲染回输入图像。初始化阶段使用闪光感知点云，让传输层高斯先落在更可信的位置；输出则是可分别新视角渲染的传输场景和反射场景。相对 Deferred Reflection，它不是给同一表面增加高光，而是把玻璃叠加的两个 3D 场景拆开。
\textbf{方法信息：} 代码：\url{https://github.com/mingyangx/flash-splat}。
\textbf{设计动机与边界：} 相对原始 3DGS 以独立高斯和图像重建误差为主的设定，Flash-Splat 把问题落在以下可实现的改动上：Flash-Splat的核心创新是将闪光/无闪光反射分离从2D配对扩展到3D非配对：(1)闪光原理：闪光增加传输场景亮度而不显著影响反射场景，通过闪光/无闪光图像差分可获得无反射传输场景；(2)非配对扩展：不再要求配对采集，先采集多视角闪光图像序列，再采集多视角无闪光图像序列，两个序列无需对齐；(3)3DGS建模：扩展3DGS分别建模传输场景和反射场景，使用闪光先验损失引导分离；(4)闪光感知点云初始化：利用闪光信息初始化3D高斯位置。方法大幅简化了图像采集过程，在真实世界实验中准确重建了传输和反射的3D场景。其设计目的不是单独增加一个模块，而是用上述表示、监督或渲染约束消除原流程在目标任务中的具体失配。提出Flash-Splat，首次将闪光/无闪光反射分离扩展到3D，去除配对图像采集的严格要求，利用3DGS的新视角合成能力实现非配对测量下的高质量3D反射分离，大幅超越现有3D反射分离方法。适用时还应保留论文给出的边界：需要闪光灯设备；对极强反射场景分离仍有挑战；假设传输场景近似漫反射
% METHOD: 2412.19282
\paragraph{Reflective Gaussian Splatting（相互反射建模）（2024.12）}
Reflective Gaussian Splatting~\cite{2412_19282} 又称 Ref-Gaussian，目标是同时重建高反射物体和物体间相互反射。它使用 2D 高斯基元增强表面几何，并采用物理延迟渲染。

管线先用 2D Gaussian surfel 替代普通 3D ellipsoid，使法线和局部表面更稳定；随后通过 $\alpha$-混合聚合 albedo、metallic、roughness 和 normal 等像素级材质属性，再在屏幕空间执行简化 Disney BRDF。镜面环境光仍采用 split-sum 近似，把环境积分分解为预滤波环境贴图和 BRDF LUT，避免每个像素做昂贵蒙特卡洛采样。

Ref-Gaussian 的增量在 inter-reflection。高反射物体之间会互相映出对方，若只查询远场环境贴图，就会把近场反射错误解释为材质或虚假几何。该方法利用高斯 splatting 的可见性推理估计高斯之间的相互反射贡献，并配合 material-aware normal propagation：当某些高斯表现出高 metallic、低 roughness 的强反射特征时，训练中放大其尺度和影响范围，让邻域高斯通过共享像素获得更可靠的法线和材质梯度。相对 Deferred Reflection，它不只是给像素加一层环境反射，而是把 2DGS 表面、PBR 材质和高斯级近场反射放进同一条优化管线。
\textbf{设计动机与边界：} 相对原始 3DGS 以独立高斯和图像重建误差为主的设定，Reflective Gaussian Splatting 把问题落在以下可实现的改动上：Ref-Gaussian的核心设计：(1) 使用2D高斯基元替代3D高斯以增强几何建模能力；(2) 基于物理的延迟渲染：先通过alpha-blending聚合像素级材料属性（albedo、metallic、roughness、normal），再在像素级执行简化版Disney BRDF渲染方程，采用split-sum近似将环境光照积分分解为预计算的LUT和卷积，避免蒙特卡洛采样；(3) 高斯级相互反射：首次在3DGS中实现相互反射，通过高斯泼溅的可见性推理替代网格光线追踪；(4) 材料感知法线传播：增大高metallic低roughness高斯基元的尺度以增强几何建模；(5) 初始高斯级着色阶段提供良好的几何初始化。其设计目的不是单独增加一个模块，而是用上述表示、监督或渲染约束消除原流程在目标任务中的具体失配。提出Ref-Gaussian，首个在3DGS框架中同时实现实时高质量反射物体重建和相互反射建模的方法。通过基于物理的延迟渲染和split-sum近似避免昂贵的蒙特卡洛采样，引入2D高斯基元和材料感知法线传播增强几何建模，在反射和非反射场景上均达到SOTA。适用时还应保留论文给出的边界：2D高斯基元在极端视角下可能出现退化；split-sum近似在复杂光照下可能引入误差；相互反射的计算开销随高斯数量增长。
\textbf{方法信息：} 代码：\url{https://github.com/fudan-zvg/ref-gaussian}；GitHub 关注度较高，约 118 stars；引用很多，Semantic Scholar 约 40 次。



% METHOD: 2506.13110
\paragraph{GS-2DGS（基础模型几何监督的反射重建）（2025.06）}
GS-2DGS~\cite{2506_13110} 以 2DGS 为基础处理反射物体。反射表面上，传统多视图几何约束弱，3DGS 的法线估计又不稳定；GS-2DGS 的改变是从基础模型引入深度和法线监督。

2D 高斯面元的法线可由切向量叉积直接得到，方法再从 DepthPro、GeoWizard/Marigold 一类基础模型获得单目深度和法线估计，通过法线一致性和深度对齐损失约束训练。外观部分沿用反射重建常见的物理分解，把材质属性和环境光照分开，使模型不仅能做新视角合成，也能输出更合理的重光照结果。相对 SpecGaussian 的潜特征外观补偿，GS-2DGS 更相信大模型几何先验对反射表面的稳定作用；相对 Ref-Gaussian，它的重点不是相互反射，而是用外部几何先验把反射表面拉回可信几何。
\textbf{设计动机与边界：} 相对原始 3DGS 以独立高斯和图像重建误差为主的设定，GS-2DGS 把问题落在以下可实现的改动上：GS-2DGS基于2DGS（面元表示，几何性能优于3DGS）构建。核心设计：(1)利用2D高斯原语的明确法线方向（切向量叉积）直接渲染法线图，避免3DGS中通过最小尺度轴估计法线的不准确性；(2)从基础模型（DepthPro/GeoWizard等）获取单目深度和法线估计作为几何约束，这些模型基于大规模训练数据对反射表面具有鲁棒性；(3)通过法线一致性损失、深度对齐损失等将基础模型先验融入2DGS训练；(4)结合物理分解：将外观分解为材质属性（反照率、粗糙度、金属度）和环境光照，支持重光照。整体方法在保持高斯泼溅的快速渲染优势同时，实现反射物体的高质量重建。为了消除原始 3DGS 在该目标任务中的具体失配，方法把上述表示、监督或渲染约束作为可学习和可执行的改动，而不是单独堆叠模块。提出结合2DGS快速渲染与基础模型几何监督的反射物体重建方法，通过深度和法线约束以及物理分解实现反射物体的精确重建和重光照，比SDF方法快一个数量级。适用时还应保留论文给出的边界：依赖基础模型的几何估计质量；物理分解的精度受限于2DGS的表示能力；对极端反射（如镜面）的处理仍有提升空间。
\textbf{方法信息：} Meta FAIR 参与了这篇论文；论文给出的仓库链接为 \url{https://github.com/hirotong/gs2dgs}，但当前只确认到 README 级别内容，完整代码仍需持续关注。



% METHOD: 2602.19916
\paragraph{Augmented Radiance Field（视角依赖增强核）（2026.02）}
Augmented Radiance Field~\cite{2602_19916} 提出增强高斯核来补偿标准 SH 对高频视角相关效果的不足。高阶 SH 内存增长快，低阶 SH 又难以捕捉尖锐高光；方法让某些高斯拥有视角依赖不透明度，显式模拟在特定方向变亮或变暗。

增强核的 opacity 由局部视角方向和可学习方向参数决定。训练先从 2DGS 初始化并得到标准渲染，再根据误差图在高误差区域自适应插入增强核，专门拟合标准 splatting 表达不了的高光和反射。相对 Latent-SpecGS 的神经解码器，增强核更轻量，也可作为后处理接入已有 splatting 框架。
\textbf{方法信息：} Tencent 参与了这篇论文；项目页显示代码计划发布，当前尚未确认完整公开代码。
\textbf{设计动机与边界：} 相对原始 3DGS 以独立高斯和图像重建误差为主的设定，Augmented Radiance Field 把问题落在以下可实现的改动上：两个核心创新：(1)增强高斯核：标准3DGS使用球谐函数(SH)编码视角相关颜色，但低阶SH无法捕捉尖锐镜面高光，高阶SH指数增长内存。增强高斯核引入视角依赖不透明度alpha(v) = sigmoid(dot(v\_local, n\_alpha) + b\_alpha)，其中v\_local是局部视角方向，n\_alpha和b\_alpha是可学习参数。这使得高斯在特定视角方向变亮或变暗，显式模拟镜面反射。(2)误差驱动补偿策略：从2DGS初始化开始（更好的表面表示），训练标准3DGS后，计算渲染误差图，在高误差区域自适应插入增强高斯核进行优化。增强核专门捕获标准SH无法表示的高频视角相关效果。最终产生增强辐射场，可作为后处理增强集成到任何splatting框架。其设计目的不是单独增加一个模块，而是用上述表示、监督或渲染约束消除原流程在目标任务中的具体失配。提出增强高斯核，通过视角依赖不透明度显式建模镜面效果。引入误差驱动补偿策略改善现有3DGS场景渲染质量。方法从2DGS初始化开始，自适应插入和优化增强高斯核，产生增强辐射场。在参数效率方面优于SOTA NeRF方法。适用时还应保留论文给出的边界：增强高斯核的视角依赖不透明度模型相对简单；误差驱动插入需要两阶段训练；增加的高斯数量需要额外存储。

\section{重光照}

重光照方法的目标是给同一几何和材质换新光照。与完整逆渲染相比，有的方法追求物理分解，有的方法只需要足够可信的光照迁移或 probe 烘焙。关键问题是光照传输：高斯如何响应新环境光、点光源或近场室内灯。
% METHOD: 2408.04474
\paragraph{LumiGauss（户外 2DGS 光照传输）（2024.08）}
LumiGauss~\cite{2408_04474} 继承 2DGS 的 surfel 几何来做户外重光照。它的出发点是：重光照不仅需要颜色，还需要稳定表面、法线、反照率和光照传输；普通 3DGS 椭球没有明确局部平面，容易把阴影和 albedo 混在一起。LumiGauss 因此把每个 2D Gaussian 看作可光照的 surfel，并为其学习 albedo、normal 和 transfer function。

光照部分用环境贴图的球谐系数表示，每张训练图像还有一个 lighting embedding，经 MLP 预测对应环境光。无阴影版本近似为 diffuse albedo 与 clamped cosine lighting；阴影版本进一步为每个 surfel 学习球谐 transfer function，表示从环境各方向到该点的可见性、阴影和部分互反射。由于环境光和 transfer 都在 SH 域内，渲染方程可用系数点积高效计算，这也是它能对接 precomputed radiance transfer 和游戏引擎的原因。相对 BiGS，LumiGauss 更依赖清晰表面法线，适合户外建筑和地标；BiGS 则更偏法线不明确的物体材料。
\textbf{实现细节：} LumiGauss（2024.08） （Relightable Gaussian Splatting in the Wild）基于 2DGS 为每个高斯建立可学习的辐射传输函数。传统做法用解析 BRDF 着色，需要估计材质参数。LumiGauss 的核心创新是直接为每个高斯参数化一组传输函数 SH 系数 $\mathbf{d}_k \in \mathbb{R}^{(n+1)^{}2}$，绕过显式 BRDF 估计。每张训练图像关联一个可学习嵌入 $e_c$，通过 MLP 预测环境光 SH 系数：$\mathbf{l}_c = \phi(e_c | W)$（$n=2$，即二阶 SH 的 9 个系数）。方法提供两种模型：无阴影模型 $c_k = \rho_k \mathbf{l}_k^T M(\mathbf{l}_c) \mathbf{l}_k$（$M$ 是从环境光 SH 推导的 $4 \times 4$ 矩阵，Ramamoorthi 闭合形式解）；有阴影模型 $c_k = \rho_k \sum_{i=1}^{(n+1)^{}2} l_c^i d_k^i$，直接学习包含阴影和互反射的传输函数。
\textbf{设计动机与边界：} 相对原始 3DGS 以独立高斯和图像重建误差为主的设定，LumiGauss 把问题落在以下可实现的改动上：使用2DGS（2D高斯溅射）替代3DGS，因为2D surfels提供平坦表面更适合重新光照。为每个2D高斯学习法线、albedo和传输函数属性。光照建模：环境光使用球谐(SH)参数化，传输函数也使用SH表示，利用SH正交性将渲染方程简化为SH系数的点积。提出非阴影和阴影两种模型：非阴影模型使用漫反射点积光照；阴影模型学习SH参数化的传输函数，可建模阴影和物体间光反射。利用SH正交性高效评估渲染方程。为了消除原始 3DGS 在该目标任务中的具体失配，方法把上述表示、监督或渲染约束作为可学习和可执行的改动，而不是单独堆叠模块。提出基于2DGS的户外场景重建和重新光照方法LumiGauss，通过球谐表示光照和传输函数实现高质量场景重建和新环境地图下的真实光照合成，支持与游戏引擎集成和快速预计算辐射度传输，在NeRF-OSR数据集上超越基线方法。适用时还应保留论文给出的边界：依赖2D surfels的质量；球谐阶数限制高频光照细节；户外场景假设环境光无限远
\textbf{端到端流程：} LumiGauss 的输入是经多视图重建得到的 2D surfels 与环境光观测；首先为每个 surfel 优化法线、albedo 和球谐传输系数，同时以球谐参数化环境光。渲染阶段利用球谐正交性把光照积分变成系数点积：非阴影路径使用漫反射项，阴影路径额外评估学习到的传输函数。训练以重建颜色为监督并联合更新材质、传输和几何属性；输出是可在新环境光下重渲染的二维高斯表面，而不是只拟合原始视角颜色。
\textbf{方法信息：} Microsoft 参与了这篇论文；代码：\url{https://github.com/joaxkal/lumigauss}。

一个需要持续关注的候选是 TranSplat~\cite{2503_22676}。它用 SH 域辐射传输做即时物体重光照，思路是根据源/目标环境光的逐法线辐照度比率调制 2DGS surfel 的外观系数，并用镜面感知的双路径 SH 传输处理高光。当前项目页 \url{https://tonyyu0822.github.io/transplat/} 显示代码仍为 Code Coming Soon；此前检索到的 \url{https://github.com/xingyoujun/transplat} 属于另一篇稀疏视角泛化 3DGS 论文，不作为本方法的公开代码证据。因此本书暂不把它作为已验证可信方法展开，后续代码发布或引用增长后应重新检查。



% METHOD: 2408.13370
\paragraph{BiGS（双向高斯基元）（2024.08）}
BiGS~\cite{2408_13370} 解决 Gaussian 表示缺少明确表面法线的问题。很多重光照方法依赖 BRDF 和法线，难以处理毛绒、半透明或体积材质；BiGS 使用双向球谐函数表示光照和视角依赖散射函数，避免显式表面模型。

它先训练普通 Gaussian splat 模型得到几何和可见性，再在 One-Light-At-a-Time 数据上学习每个高斯的 lighting-dependent appearance。核心参数不是法线、roughness 和 metallic，而是一组双向球谐系数：一端描述入射光方向，另一端描述观察方向，二者共同决定该高斯在某个光源和视角下的散射响应。外观被拆成 diffuse scattering 和 directional scattering 两部分，并用非负性、能量约束等正则防止方向散射任意放大。推理时，点光源、方向光或环境贴图先编码到 SH 域，再与每个高斯的双向系数求值，最后进入 splatting。相对 LumiGauss 的 surfel 法线和 transfer function，BiGS 更适合毛绒、半透明、玉石这类边界模糊或体积感强的材质。
\textbf{实现细节：} 前述方法（PRTGS、RNG、F-RNG）分别用解析 BRDF 或神经网络建模每个高斯的反射行为，但都假设散射函数可以拆分为独立的入射和出射项。BiGS（2024.08） （Bidirectional Gaussian Primitives）修改了 3DGS 的颜色编码环节——用双向球谐（Bidirectional SH）同时参数化入射和出射方向的联合散射函数——高斯的几何属性保持不变。双向球谐与互易性。传统做法将高斯颜色参数化为视角相关的球谐系数，无法响应光照变化。BiGS 将每个高斯的颜色分解为直接光照 $L_{\mathrm{direct}}(\omega_o)$ 和间接光照 $L_{\mathrm{ind}}$，散射函数进一步分解为漫反射项 $\mathbf{d} \in \mathbb{R}^{}3$（视角无关的 albedo+环境遮蔽）和方向散射项 $s(\omega_i, \omega_o)$。关键创新是用双向球谐 $s(\omega_i, \omega_o) = \sum_{i=1}^{n} \sum_{j=1}^{n} c_{ij} y_i(\omega_i) y_j(\omega_o)$ 表示方向散射——直觉上，这是一张二维"查找表"：一个维度是光从哪来，另一个维度是你从哪看，表中的值告诉你表面有多亮。BiGS 并显式强制互易性 $c_{ij} = c_{ji}$（物理上散射函数必须满足 $f(\omega_i, \omega_o) = f(\omega_o, \omega_i)$——交换光源和观察者位置，看到的亮度不变），将参数从 $n^{}2$ 减少到 $n(n+1)/2$。使用 25 个 SH 基函数，$s$ 仅需 $325 \times 3 = 975$ 个参数。高效重光照。重光照时，先对 $\omega_i$ 求和得到中间系数 $c_j = \sum_{i} c_{ij} y_i(\omega_i)$，再对新视角 $\omega_o$ 求值——将球面积分转化为两次球谐求值，复杂度 $O(n^{}2)$。
\textbf{方法信息：} 代码：\url{https://github.com/desmondlzy/bigs}。
\textbf{设计动机与边界：} 相对原始 3DGS 以独立高斯和图像重建误差为主的设定，BiGS 把问题落在以下可实现的改动上：BiGS将光内在分解集成到Gaussian Splatting框架中。核心设计：(1)外观模型采用双向球谐函数(BSH)表示光照和视角依赖的散射函数，不依赖表面法线相关的反射函数，因此兼容Gaussian这种法线未定义的体积表示；(2)散射分解为漫散射和方向散射分量，作为光照依赖辐射的内在分解；(3)物理启发的正则化项确保优化过程的稳定性；(4)使用One-Light-At-a-Time(OLAT)数据作为输入训练。模型能够处理近场和环境光照，渲染毛绒、半透明等体积材质效果。其设计目的不是单独增加一个模块，而是用上述表示、监督或渲染约束消除原流程在目标任务中的具体失配。提出双向Gaussian原语(Bidirectional Gaussian Primitives)，首次在Gaussian Splatting框架中统一表面和体积材质的重光照，使用双向球谐函数避免了对表面法线的依赖，实现实时动态光照渲染。适用时还应保留论文给出的边界：需要OLAT数据作为输入，限制了在一般场景中的适用性；双向球谐函数的阶数影响表示能力和计算量。
% METHOD: 2409.19702
\paragraph{RNG（Relightable Neural Gaussians）（2024.09）}
RNG~\cite{2409_19702} 同样避免强表面着色假设。每个高斯携带潜在特征，神经解码器根据位置、观察方向和光照方向输出辐射率，从而支持硬表面和软边界物体的重光照。

方法引入阴影线索和深度精化网络来处理点光源产生的锐利阴影，并使用混合前向-延迟拟合策略：前向优化阶段保证几何质量，延迟优化阶段提升外观细节。训练数据来自移动点光源下的多视图观测，让模型看到多个视角/光照组合以减少分解歧义。相对 BiGS 的解析双向球谐，RNG 更依赖神经解码器；相对 F-RNG，它仍是每个对象单独优化，因此质量更容易贴合单场景，泛化和速度则不如前馈路线。
\textbf{实现细节：} Gaussian in the Wild（GS-W）（2024.03） （3DGS for Unconstrained Image Collections）首次将 3DGS 应用于互联网照片集合，核心挑战是同时处理光照变化和瞬态物体。标准 3DGS 使用球谐系数建模颜色，但 GS-W 放弃了这一方案，转而为每个高斯引入外观特征 $af_i$，将其分解为固有外观 $sf_i$（可学习属性，建模材质本身的颜色）和动态外观 $df_i$（从参考图像提取的特征，建模光照引起的变化），两者通过融合 MLP 生成最终外观特征 $af_i = M_f(sf_i, df_i, X_i)$，再由颜色解码 MLP 得到高斯颜色 $c_i = M_c(af_i, d)$。动态外观特征的提取采用自适应采样策略：从参考图像通过 UNet-ResNet 提取 1 张投影特征图 $F^P$ 和 $K$ 张附加特征图 $(F^{}1, ..., F^K)$，每个高斯拥有 $K$ 个可学习采样坐标 $(sc_i^{}1, ..., sc_i^K)$，自适应决定在哪些位置采样特征 $df_i = f_i^P \oplus f_i^{}1 \oplus ... \oplus f_i^K$，采样坐标通过正则化 $\mathcal{L}_{sc} = \frac{1}{N}\sum_i \max(0, |sc_i| - 1)$ 约束在有效范围内。
\textbf{设计动机与边界：} 相对原始 3DGS 以独立高斯和图像重建误差为主的设定，RNG 把问题落在以下可实现的改动上：RNG方法：(1)可重光照神经高斯：每个高斯点携带潜在特征向量，通过神经高斯解码器phi解码为视角和光照相关的辐射率：R(x,o\_v,l\_i)=phi(x|o\_v,l\_i)。避免解析着色模型假设，适用于不符合简单模型的材质；(2)阴影线索：点光源产生锐利阴影，对神经网络挑战大。引入阴影线索辅助阴影学习，结合深度精化网络提高阴影准确性；(3)混合前向-延迟拟合策略：平衡几何质量和外观质量，前向优化关注几何，延迟优化关注外观细节；(4)训练设置：使用移动点光源拍摄的多视图数据，观察多种视角/光照组合减少分解歧义。其设计目的不是单独增加一个模块，而是用上述表示、监督或渲染约束消除原流程在目标任务中的具体失配。提出RNG，首个不依赖表面着色模型和表面约束的3DGS可重光照框架。使用神经高斯隐式建模辐射率，支持硬表面和软边界物体。引入阴影线索+深度精化网络改善阴影质量，混合前向-延迟拟合策略平衡几何与外观。训练1.3小时，渲染60FPS。适用时还应保留论文给出的边界：需要移动点光源的专业采集设置；阴影质量仍受限于深度估计精度；复杂全局光照效果难以捕捉。
\textbf{方法信息：} Adobe Research 参与了这篇论文；代码：\url{https://github.com/sssssy/RNG_release}。
% METHOD: 2410.06231
\paragraph{RelitLRM（可重光照大重建模型）（2024.10）}
RelitLRM~\cite{2410_06231} 是稀疏视角前馈重光照路线。传统逆渲染需要密集捕获和慢优化；RelitLRM 从 4 到 8 张未知光照输入图像，在约秒级生成可重光照 3DGS。

模型基于 GS-LRM。几何重建器用 transformer 从输入图像和 Plucker rays 生成高斯位置、协方差和不透明度；外观生成器不做确定性预测，而是用 relit-view 扩散模型在高斯外观参数空间去噪，条件包括目标光照和几何特征。关键在于扩散模型用 $x_0$ 预测目标，并通过 3D 表示瓶颈维持多视角一致：扩散模型可以生成高光、阴影和互反射这类多模态外观，但最终仍要落回同一个 3DGS 表示，避免每个视角各画各的。相对 F-RNG，它是更基础的前馈重光照模型；F-RNG 则试图在不重训 LRM 的情况下补可重光照能力。
\textbf{实现细节：} Unifying Color and Lightness Correction（2026.02） （View-Adaptive Curve Adjustment）解决不同训练图像之间曝光和白平衡不一致的问题。标准 3DGS 假设所有图像在统一颜色空间中，但实际上每张照片的相机 ISP 处理链不同。该方法为每张训练图像学习一组色调曲线参数，通过视角自适应的曲线调整将不同曝光和白平衡的图像映射到统一颜色空间，使 3DGS 的光度损失在一致的色彩空间中计算。
\textbf{方法信息：} Adobe Research 参与了这篇论文；项目页已公开，当前未确认完整代码仓库。
\textbf{设计动机与边界：} 相对原始 3DGS 以独立高斯和图像重建误差为主的设定，RelitLRM 把问题落在以下可实现的改动上：RelitLRM基于GS-LRM构建，包含两个核心组件：(1)几何重建器：使用与GS-LRM相同的transformer架构，从输入图像和Plucker射线生成3DGS几何参数（位置、协方差、不透明度）；(2)可重光照外观生成器：用relit-view扩散模型替代GS-LRM的确定性外观预测。扩散模型以目标光照条件和预测的几何特征为条件，在3DGS外观参数空间进行去噪，生成重光照后的球谐系数。关键创新在于外观生成器使用x0预测目标而非噪声预测，并通过3D表示瓶颈确保多视角一致性。整个模型在合成多视角多光照数据上端到端训练。其设计目的不是单独增加一个模块，而是用上述表示、监督或渲染约束消除原流程在目标任务中的具体失配。提出RelitLRM，首个基于前馈transformer的可重光照大重建模型，通过扩散模型捕获重光照外观的多模态分布（如镜面高光和阴影），在约1秒内从4-8张稀疏图像生成可重光照的3DGS表示适用时还应保留论文给出的边界：训练依赖大规模合成数据，在真实世界复杂材质上的泛化能力有待验证；扩散模型推理增加了一定计算开销
% METHOD: 2507.02257
\paragraph{GBake（高斯场景烘焙为反射探针）（2025.07）}
GBake~\cite{2507_02257} 不直接求 BRDF，而是把 3DGS 场景烘焙成游戏引擎可用的反射探针。高斯场景把几何和光照编码在外观里，传统 mesh 插入后无法正确接收场景反射；GBake 通过一组局部环境贴图弥合两种管线。

具体做法是在训练好的高斯场景包围盒内放置规则 3D probe grid；每个网格点从六个方向渲染 cubemap，得到该位置看到的局部高斯环境。输出是一组 cross-layout cubemap PNG 和记录 probe 三维位置的 JSON，Unity 等引擎可按传统 reflection probe 方式为插入 mesh 查询附近环境贴图。实现上提供 threedgrut/3D-GRT 光线追踪后端和 gsplat 后端，输入通常是训练好的 \texttt{.ply} 高斯模型。相对 LuxRemix 的“拆光源并交互调灯”，GBake 更像工程桥接：它不恢复可编辑 BRDF，也不求全局照明，只把已经烘焙在 3DGS 外观中的局部环境转成传统图形管线能消费的反射资源。
\textbf{方法信息：} 代码：\url{https://github.com/gbake-official/gbake}。
\textbf{设计动机与边界：} 相对原始 3DGS 以独立高斯和图像重建误差为主的设定，Gbake 把问题落在以下可实现的改动上：GBake的核心思想是将3DGS场景烘焙为规则网格分布的环境贴图（反射探针）。由于3DGS将光照和几何联合编码为外观，直接插入mesh无法正确重光照，GBake通过创建场景内的反射探针网格来编码异构光照。具体步骤包括：(1) 在场景中建立规则3D网格；(2) 在每个网格点使用3DGS渲染全景环境贴图；(3) 将这些探针导入Unity引擎，利用反射映射技术为插入的mesh对象提供局部光照信息。该方法避免了需要显式光源和表面几何的困难，通过多个环境贴图编码场景内的光照异构性。其设计目的不是单独增加一个模块，而是用上述表示、监督或渲染约束消除原流程在目标任务中的具体失配。提出GBake工具，将3DGS场景烘焙为反射探针（环境贴图）网格，使传统mesh对象在Unity游戏引擎中能够获得真实的反射映射效果，有效弥合了外观优化的高斯表示与传统mesh渲染管线之间的差距。适用时还应保留论文给出的边界：反射探针的网格密度需要在精度和存储之间权衡；对于高频光照变化（如聚光灯附近）可能需要更密的探针网格；仅解决了反射映射问题，未处理阴影和全局光照。
% METHOD: 2601.15283
\paragraph{LuxRemix（室内多光源拆分与编辑）（2026.01）}
LuxRemix~\cite{2601_15283} 关注室内交互式光照编辑。室内场景不是一个远场环境光能解释的：台灯、顶灯、窗光和局部反射都有空间变化。LuxRemix 的改变是用生成式图像光照分解模型把多视角捕获分解为可独立控制的光源贡献。

流程分三步。第一步先训练或使用单图像光照分解模型，把每个输入视图拆成 one-light-at-a-time 光源贡献和 ambient 分量，用户可以分别控制光源开关、色度和强度。第二步做 multi-view lighting harmonization：同一个灯在不同视角中的分解结果必须对齐，否则从一个视角关灯、换到另一个视角又会变成不同效果。第三步把协调后的光照分解写入可重光照 3DGS，使高斯场不仅能渲染原始外观，还能按指定光源组合实时合成新视角。相对 LumiGauss 的远场环境光 SH，LuxRemix 更偏近场、多光源和室内交互控制；相对 GBake 的反射探针烘焙，它直接面向“拆灯并重混”的编辑任务。
\textbf{实现细节：} LuxRemix（2026.01） （Lighting Decomposition and Remixing）将室内场景光照分解为单独可控的光源贡献。传统做法需要物理上逐一开关光源采集 OLAT 数据，实际操作困难。
\textbf{实现补充：} 实现上还要核对输入、输出、参数、特征、初始化、训练损失、正则、渲染或推理接口，以及优化中的剪枝、分裂、采样、编码或解码步骤；这些抓手决定该方法到底改了 3DGS 管线的哪一环。
\textbf{方法信息：} Meta Reality Labs 参与了这篇论文；Meta 公开了 LuxRemix 数据集仓库 \url{https://github.com/facebookresearch/luxremix-dataset}，但该仓库不是方法完整实现。
\textbf{设计动机与边界：} 相对原始 3DGS 以独立高斯和图像重建误差为主的设定，LuxRemix 把问题落在以下可实现的改动上：方法分三个阶段：(1)使用生成式图像光照分解模型将室内场景照明分解为独立光源贡献，支持对每个光源的状态(开/关)、色度和强度进行独立操控；(2)多视角光照协调模块确保分解结果在所有视角间一致传播；(3)训练可重照明的3DGS表示，支持从任意视角实时交互操控各光源。不同于将光照视为全局属性或聚焦远距离照明的方法，本方法建模复杂室内环境中多近场光源的空间变化全局照明。其设计目的不是单独增加一个模块，而是用上述表示、监督或渲染约束消除原流程在目标任务中的具体失配。提出LuxRemix方法，从单次多视角捕获实现室内场景交互式光照编辑。核心包括：(1)生成式图像光照分解模型将复杂室内照明分解为可独立控制的独立光源；(2)多视角光照协调确保跨视图一致性；(3)集成到可重照明的3DGS表示中，实现实时交互控制各光源的开/关、色度和强度。适用时还应保留论文给出的边界：依赖单次多视角捕获质量；对极端光照变化（如完全关灯后的补光）效果可能受限；复杂多光源交互的物理准确性需进一步验证。
% METHOD: 2605.25975
\paragraph{F-RNG（前馈可重光照神经高斯）（2026.05）}
F-RNG~\cite{2605_25975} 把 RNG 的 per-scene 优化路线前馈化。现有可重光照方法往往需要密集视角和每场景优化；前馈 LRM 虽快，却常把光照烘焙到外观里。F-RNG 的目标是从稀疏视角直接生成可重光照高斯资产。

它分三步。先检测显著补丁，在 LRM 潜空间插值生成精细几何 token，并用专门 detokenizer 解码为高斯，IDM 先验提供法线和反照率引导；这一步解决的是前馈 LRM 常见的几何细节不足。然后 MaterialFormer 从 RelitLRM 的辐射输出中蒸馏可重光照外观表示，如反照率、金属度和粗糙度，并用光照独立性正则约束，避免把输入光照继续烘焙进材质。最后通用神经渲染器接收高斯几何、外观参数和目标环境光，用小型 MLP 输出最终颜色。相对 RelitLRM，它不重新训练底层 LRM，而是训练额外小网络；相对 RNG，它从稀疏视图一次前馈生成资产，不再为每个对象做长时间优化。
\textbf{实现细节：} RNG 对每个场景单独训练解码器 MLP，耗时数分钟，无法快速处理新物体。F-RNG（2026.05） （Feed-Forward Relightable Neural Gaussians）修改了 RNG 的优化环节——用前馈 Transformer 替代逐场景 MLP 优化——输入的多视角图像表示保持不变。F-RNG 的核心改进是利用 RelitLRM 的前馈 Transformer 预测几何和外观 token，再通过 MaterialFormer 提取可重光照的神经材质表示。具体分三步。(1) 潜空间插值精细化几何。检测纹理复杂度最高的 top-$K$ 图像块，在 RelitLRM 的 latent 空间中插值生成亚像素级精细高斯，通过专用 de-tokenizer 解码。(2) 先验引导的可重光照外观蒸馏。训练 MaterialFormer Transformer，将 RelitLRM 的几何/外观 token 和扩散先验 token 融合为材质 token $T_i^{\mathrm{mat}}$，并施加光照无关正则化：对同一物体输入不同环境光 $E_0, E_1$，强制材质 token 一致： 其中 $\mathcal{C}$ 为余弦相似度，$\mathcal{L}_{\mathrm{KLD}}$ 为 KL 散度。直觉是：材质是物体的固有属性，不应该因为换了光照就改变。(3) 通用神经渲染器。不同于 RNG 的逐场景 MLP，F-RNG 训练一个跨所有物体共享的通用外观解码器 $\rho(G_i^{\mathrm{mat}}, \omega_o, \omega_i) = \Phi(G_i^{\mathrm{mat}}, \omega_o, \omega_i)$，支持从皮草到金属的广泛材质类型。
\textbf{方法信息：} NVIDIA 参与了这篇论文；当前未确认完整公开代码。



\section{HDR与光照变化}

HDR 和光照变化路线处理的是相机与采集条件，而不是材质本身。野外照片集存在不同时间、天气、曝光和瞬态遮挡；HDR/RAW 方法则要求在极亮和极暗区域都保留辐射信息。这里的锚点是 GS-W、HDR-GS 和 LE3D：GS-W 拆外观变化，HDR-GS 拆动态范围和相机响应，LE3D 则把 3DGS 推到低光 RAW 线性空间，重点解决低信噪比、远景 SfM 稀疏和 RAW 颜色表达不足。


% METHOD: 2405.15125
\paragraph{HDR-GS（双动态范围高斯）（2024.05）}
HDR-GS~\cite{2405_15125} 是 HDR 3DGS 的代表性锚点。原始 3DGS 假设训练图像已经是统一曝光的 LDR 颜色，遇到多曝光输入时，SH 颜色会被迫同时解释曝光变化和真实辐射，容易不收敛或产生颜色伪影。HDR-GS 的改变是把每个高斯的颜色拆成 HDR radiance 和曝光后的 LDR color，并让二者通过可学习相机响应相连。

具体来说，每个高斯仍有位置、协方差和 opacity，但颜色侧维护 HDR SH 系数；给定用户输入或图像元数据中的曝光时间 $t$，三个 RGB 通道的 MLP 在对数域学习 inverse camera response function，把 $\log$ HDR radiance 与 $\log t$ 映射成 LDR 颜色。训练时有两条并行可微 rasterization：HDR 分支直接渲染 HDR 图，与由多曝光图合成的 HDR 参考对齐；LDR 分支先经 tone mapper 得到曝光相关颜色，再渲染并与真实 LDR 图像计算 $L_1$ 和 SSIM 损失。论文还重新标定 HDR-NeRF 数据集的相机并用 SfM 点初始化高斯，使 3DGS 能稳定进入 HDR NVS 设置。相对后续 HDRSplat，HDR-GS 更像“多曝光 LDR 到 HDR/LDR 双输出”的基础框架；相对 Cinematic Gaussians，它还没有把光圈和焦距纳入相机模型。
\textbf{实现细节：} HDR-GS的核心是双动态范围（DDR）高斯点云模型。每个3D高斯包含位置、协方差、不透明度、LDR颜色和HDR颜色属性。HDR颜色直接用球谐系数建模，LDR颜色通过色调映射从HDR颜色和可控曝光时间t推导。为避免MLP训练的梯度消失问题，色调映射在离散对数域进行：将CRF的逆函数用三个独立MLP（RGB三通道）学习，先在离散对数域建模再反变换回LDR颜色。训练时，LDR和HDR颜色分别通过两个并行可微光栅化（PDR）过程渲染对应图像。LDR渲染图与真实LDR图像计算损失，HDR渲染图与合成的HDR参考图计算损失。此外，重新标定相机参数并计算高斯初始位置，为HDR 3DGS研究建立数据基础。总损失函数包含LDR和HDR的L1+SSIM损失。 提出HDR-GS框架，首次实现基于3DGS的高效HDR新视角合成。设计双动态范围（DDR）高斯点云模型，用球谐系数拟合HDR颜色并通过MLP色调映射器渲染LDR颜色。两个并行可微光栅化（PDR）过程同时重建HDR和LDR视图。相比SOTA NeRF方法在LDR和HDR NVS上分别提升3.84和1.91dB，推理速度提升1000倍，训练时间仅需6.3
\textbf{方法信息：} 代码：\url{https://github.com/caiyuanhao1998/hdr-gs}；GitHub 关注度较高，约 388 stars。



% METHOD: 2406.06216
\paragraph{LE3D（RAW 低光 HDR 高斯）（2024.06）\cite{2406_06216}}
LE3D 处理的是另一类 HDR 输入：不是多曝光 LDR 图像，而是夜间或弱光场景中的 noisy RAW 图像。RawNeRF 证明了在 RAW 线性空间中做视图合成可以保留暗部和高动态范围信息，但体渲染训练和渲染都慢；原始 3DGS 虽然快，却有三个直接障碍：低光图像让 COLMAP 难以给远景提供足够点云，基于 SH 的颜色项不适合无上界的 RAW 线性颜色，稀疏高斯在噪声监督下还容易形成不稳定的几何。LE3D 的改变是把 RawNeRF 的 RAW/HDR 监督思想迁到 3DGS，同时修改初始化、颜色参数化和结构正则。

它的流程先运行 COLMAP 得到相机与稀疏点云，然后用 Cone Scatter Initialization 在所有训练相机共同覆盖的视锥 $\mathcal{F}$ 内补充远景点。视锥由平均相机方向、最大视场角以及 COLMAP 点到虚拟视点的近远平面确定，补点后得到 $\mathcal{S}'=\mathcal{S}\cup\mathcal{S}^{\mathcal{F}}$，再用 $\mathcal{S}'$ 初始化高斯。颜色侧，LE3D 不再用 SH 直接拟合 RAW 颜色，而是为每个高斯维护颜色特征 $f_i$ 和 primitive-aware bias $b_i$，把 $f_i$ 与相机视角 $v$ 输入一个小型 Color MLP，并用指数激活输出无上界的线性颜色：
\begin{equation}
c_i=\exp\left(\mathbf{F}_{\theta}(f_i,v)+b_i\right),\qquad
b_i^{(0)}=\log\left(\operatorname{avg}\{c_{\mathrm{pix}}\}_i\right).
\end{equation}
这样初始化时每个高斯的颜色接近训练图像投影到该点的平均 RAW 像素，优化时又能表达视角相关和高动态范围颜色。损失上，LE3D 继承 RawNeRF 的加权 RAW 监督，并加入 depth distortion 与 near-far regularization，让同一条光线上的高斯更集中在合理深度范围；这对后续 refocus 尤其重要。相对 HDR-GS，LE3D 的核心不是从多曝光 LDR 学 CRF，而是在 RAW 线性空间训练 3DGS；相对 HDRSplat，它更早把低光初始化、Color MLP 和结构正则组合成一套可实时渲染的 RAW/HDR 高斯流程。
\textbf{实现细节：} 此外，Gaussian-DK 观察到训练过程中少量高斯会试图补偿无法被相机响应模块完全建模的复杂亮度变化，它们在自适应密度控制中倾向于远离场景中心、靠近相机，形成"浮游物"伪影。为此提出基于训练迭代的梯度缩放策略，根据高斯在归一化设备坐标（NDC）中的深度计算距离权重 $w(i,j)$，在训练早期（$t < t_1$）对近相机高斯施加强梯度抑制以防止分裂和克隆，中期（$t_1 \leq t < t_2$）逐步退火，后期（$t \geq t_2$）完全释放以恢复高频细节。Gaussian-DK 还贡献了一个面向暗光场景的多视角基准数据集，并展示可通过控制曝光水平合成不同亮度的"照亮"图像。LE3D（2024.06） 面向 RAW 图像的 HDR（High Dynamic Range，高动态范围，能够同时捕捉极亮和极暗细节的成像技术）视角合成，提出锥形散射初始化丰富 SfM 点云估计，用 Color MLP 替代球谐函数表示 RAW 线性色彩空间，并引入深度畸变和远近正则化改善场景结构。From Chaos to Clarity（2024.06） 揭示了噪声对 3DGS 的特殊影响：3DGS 极易产生过拟合噪声的细长高斯形状，因此提出自监督框架集成噪声提取器和噪声鲁棒重建损失。
\textbf{设计动机与边界：} 相对原始 3DGS 以独立高斯和图像重建误差为主的设定，Lighting Every Darkness with 3DGS 把问题落在以下可实现的改动上：LE3D针对3DGS直接应用于夜间RAW数据的三个核心问题提出解决方案：(1) 锥形散射初始化(CSI)：估计所有相机的位置和朝向，在预定义视锥内随机散射点来丰富COLMAP点云，覆盖远距离场景。(2) Color MLP替代SH：原始SH无法充分表示RAW线性色彩空间的HDR信息。每个高斯初始化随机颜色特征和偏置，通过轻量MLP预测颜色。(3) 深度失真正则化和近远正则化：改善场景结构精度，支持下游任务如重聚焦。支持实时HDR渲染、色调映射变化和重聚焦功能。其设计目的不是单独增加一个模块，而是用上述表示、监督或渲染约束消除原流程在目标任务中的具体失配。提出LE3D解决3DGS在夜间RAW数据上的三大问题：(1)锥形散射初始化丰富SfM点云估计；(2)用Color MLP替代球谐函数表示RAW线性色彩空间；(3)深度失真和近远正则化改善场景结构精度。相比体积渲染方法，训练时间降至1%，2K分辨率渲染速度提升4000倍。适用时还应保留论文给出的边界：对极度噪声的RAW图像仍需进一步优化；Color MLP增加少量推理开销。
\textbf{方法信息：} 代码：\url{https://github.com/Srameo/LE3D}；GitHub 关注度较高，约 234 stars；Semantic Scholar 约 30 次引用。



% METHOD: 2406.07329
\paragraph{Cinematic Gaussians（HDR 与景深）（2024.06）}
Cinematic Gaussians~\cite{2406_07329} 在 HDR-GS 的动态范围问题上再加一个摄影机自由度：真实照片往往还有不同光圈和焦距，输入并不都是 pinhole、all-in-focus 图像。它的目标是从多视角 LDR 图像重建一个 HDR 3DGS，同时允许训练后调曝光、调光圈、调焦点，生成可控景深。

流程先用 SfM 标定相机，但由于焦距、光圈等镜头参数以毫米计，方法还利用输入图像中的散焦区域估计场景绝对尺度：对图像做清晰度检测，找出 in-focus 像素，用已知 focus distance 约束 vanilla 3DGS 渲染深度的尺度。表示上，高斯仍存 HDR SH 颜色；渲染时先按薄透镜模型把投影后的高斯与高斯光圈函数做解析卷积，得到散焦后的屏幕高斯，再进入 rasterizer。最后，tone mapping 模块把 HDR 渲染映射到 LDR 观测用于训练。训练完成后，\texttt{render\_taf.py} 这类后处理入口可以改变 exposure、aperture 和 focus distance。相对 HDR-GS，它不是只控制曝光，而是把镜头景深作为可微相机模型的一部分。
\textbf{实现细节：} Cinematic Gaussians（2024.06） （Real-Time HDR Radiance Fields with Depth of Field）在 3DGS 中同时实现 HDR 渲染和物理正确的景深效果。标准 3DGS 假设针孔相机模型（所有物体都清晰），但真实镜头受光圈和焦距影响，离焦物体会产生模糊。Cinematic Gaussians 引入薄透镜模型，每个投影高斯的弥散圆（circle of confusion）半径为 $r(z) = \frac{|z - z_f|}{z} \cdot \frac{f^{}2}{2N(z_f - f)}$，其中 $z$ 为高斯深度、$z_f$ 为对焦距离、$f$ 为焦距、$N$ 为光圈 F 值。弥散圆用高斯近似，其协方差为 $\Sigma_a = \frac{(3r(z), 3r(z))^{}2}{2}$，与场景高斯的 2D 投影协方差卷积：$\tilde{\Sigma}_i = \Sigma_i + \Sigma_a$，闭式更新后的高斯核为 $\tilde{\alpha}_i(x) = \alpha_i \mathcal{N}(x; \mu_i, \tilde{\Sigma}_i) \cdot \frac{|\Sigma_i|}{|\tilde{\Sigma}_i|}$，其中分式项为能量保持归一化因子。HDR 管线将高斯球谐系数视为 HDR 辐亮度值，渲染后通过显式色调映射函数转换为 LDR：$T(E) = \mathrm{clip}(\frac{E \cdot t \cdot A}{N^{}2}, 0, 1)^{1/\gamma}$，其中 $t$ 为曝光时间、$A$ 为光圈面积、$\gamma$ 为 gamma 校正。该公式将物理镜头参数直接嵌入色调映射，使用户可以在渲染时实时调整光圈大小 $N$、曝光时间 $t$ 和对焦距离 $z_f$，获得电影级的景深和曝光控制。训练损失包含三项：$\mathcal{L} = \mathcal{L}_r + \lambda_s \mathcal{L}_s + \lambda_d \mathcal{L}_d$，其中 $\mathcal{L}_r$ 为直接光度损失、$\mathcal{L}_s$ 为归一化曝光后的光度损失（消除不同曝光时间的影响）、$\mathcal{L}_d$ 为可微模糊检测器损失（匹配输入图像的景深模糊分布），$\lambda_s = 0.25$，$\lambda_d = 1$。
\textbf{方法信息：} 代码：\url{https://github.com/Hans1984/CineGS}。
\textbf{设计动机与边界：} 相对原始 3DGS 以独立高斯和图像重建误差为主的设定，Cinematic Gaussians 把问题落在以下可实现的改动上：Cinematic Gaussians接受多视图LDR图像（不同曝光时间、光圈和焦距）作为输入。核心技术贡献：(1) 薄透镜景深模型：将景深建模为投影场景高斯与高斯光圈函数的高效卷积，基于薄透镜相机模型。卷积后的高斯通过光栅化器渲染，实现物理准确的散景效果。(2) HDR表示：3D高斯存储HDR颜色值（通过SH系数），通过色调映射模块将HDR渲染结果映射到LDR用于训练。(3) 相机标定：利用散焦图像和已知镜头参数恢复场景绝对尺度，实现物理正确的景深建模。训练完成后支持新视角合成中任意组合的曝光、光圈和焦距参数。其设计目的不是单独增加一个模块，而是用上述表示、监督或渲染约束消除原流程在目标任务中的具体失配。提出基于3DGS的轻量方法，利用多视图不同曝光、光圈和焦距的LDR图像重建HDR辐射场。通过基于薄透镜模型的高斯解析卷积实现景深效果，结合色调映射模块实现实时电影级渲染。同时处理HDR和景深，超越现有SOTA。适用时还应保留论文给出的边界：需要多视图不同曝光和焦距的输入图像；薄透镜模型无法模拟复杂光学畸变。
% METHOD: 2407.12306
\paragraph{Splatfacto-W（Nerfstudio 野外高斯实现）（2024.07）}
Splatfacto-W~\cite{2407_12306} 可以看作 GS-W 思路在 Nerfstudio 生态中的实用化。它为每个高斯分配外观特征，为每张训练图像维护外观嵌入，通过小 MLP 预测颜色调整；训练完成后把 MLP 输出烘焙为显式颜色，推理时不再查询 MLP，保持 3DGS 速度。

瞬态物体处理采用启发式残差统计，而非依赖 2D 预训练分割模型：某区域若只在部分视角异常，就被视为行人、车辆等瞬态。天空和远景由球谐背景模型单独表示，避免前景高斯承担背景光照变化。相对 GS-W，它的贡献更多在可复现框架、烘焙推理和工程稳定性。
\textbf{设计动机与边界：} 相对原始 3DGS 以独立高斯和图像重建误差为主的设定，Splatfacto-W 把问题落在以下可实现的改动上：Splatfacto-W的三个核心贡献：(1)潜在外观建模：为每个3D高斯分配外观特征向量，为每张训练图像维护外观嵌入。两者通过小型MLP预测颜色调整，训练完成后将MLP输出烘焙为显式颜色，确保推理时无需查询MLP，保持3DGS的渲染速度。这解决了3DGS显式表示难以处理光照变化的问题；(2)瞬态物体处理：提出基于启发式的方法掩码瞬态物体（如行人、车辆），不依赖2D预训练模型。通过分析渲染残差的统计特性（如某区域在部分视角异常）自动识别和掩码瞬态区域；(3)球谐背景模型：使用球谐函数建模天空和远景背景，解决3DGS背景建模不足的问题。背景模型与前景高斯分开优化，更好地表示不同光照条件下的天空颜色变化。方法集成在Nerfstudio框架中，易于使用和扩展。其设计目的不是单独增加一个模块，而是用上述表示、监督或渲染约束消除原流程在目标任务中的具体失配。实现Nerfstudio中首个处理不受控照片集的3DGS方法。通过每高斯外观特征+每图像外观嵌入的潜在外观建模、启发式瞬态物体处理和球谐背景模型，PSNR平均提升5.3dB，训练比NeRF方法快150倍，渲染速度与3DGS相当。适用时还应保留论文给出的边界：启发式瞬态检测可能在复杂场景中失效；外观建模的MLP增加训练时间；背景模型假设天空为球面可能不适用于室内
\textbf{方法信息：} 代码：\url{https://github.com/KevinXu02/splatfacto-w}；GitHub 关注度较高，约 133 stars；引用很多，Semantic Scholar 约 44 次。



% METHOD: 2407.16503
\paragraph{HDRSplat（RAW 空间 HDR 重建）（2024.07）}
HDRSplat~\cite{2407_16503} 直接从 14 位线性 RAW 图像训练 3DGS。夜间和低光场景中，ISP 后的 LDR 图像会丢失暗部和亮部信息，COLMAP 初始化也常不均匀；HDRSplat 在 RAW/HDR 空间设计损失和光栅化调优。

方法先在 Bayer RAW 空间做深度去噪，并把去噪后的线性 RAW 作为监督，而不是先经过相机 ISP 变成 8-bit JPEG。损失函数在 14 位线性空间中设计，暗区不会因为数值小而被忽略，近饱和亮区也不会简单截断；同时用 stop-gradient 的尺度化 $L_1$/DSSIM 形式稳定噪声区域训练。另一方面，HDRSplat 调整 3DGS 的 densification 和 rasterization 超参数，减少低纹理、低照度区域对 COLMAP 点云的依赖，并让模型用更少高斯覆盖夜景结构。训练后，线性 HDR 表示可接自定义 tone mapping、曝光调节、合成散焦和深度提取。相对 HDR-GS，它更强调 RAW 传感器数据和夜间 HDR 重建；相对 Cinematic Gaussians，它不把光圈/焦距作为主要建模对象。
\textbf{方法信息：} 代码：\url{https://github.com/shreyesss/hdrsplat}。
\textbf{设计动机与边界：} 相对原始 3DGS 以独立高斯和图像重建误差为主的设定，HDRSplat 把问题落在以下可实现的改动上：HDRSplat的核心贡献：(1) HDR空间专用损失：设计线性HDR空间损失函数，能有效同时从噪声暗区和近饱和亮区提取场景信息，并通过设计处理视角相关颜色而不增加球谐度数。(2) 光栅化调优：通过仔细调整光栅化参数，隐式克服3DGS对COLMAP SfM点云初始化质量的严重依赖——这对低纹理、高景深、低光照区域尤为重要，因为这些区域的COLMAP点云通常密度不均。(3) 深度Bayer空间去噪：对原始RAW图像进行去噪处理作为监督信号。整体流程直接在14位线性RAW空间训练3DGS，跳过有损的ISP后处理，保留完整动态范围。训练完成后支持合成散焦、稠密深度图提取、后控曝光/色调映射/视角等应用。其设计目的不是单独增加一个模块，而是用上述表示、监督或渲染约束消除原流程在目标任务中的具体失配。提出HDRSplat，首个基于3DGS的14位线性HDR 3D场景重建方法。设计HDR空间专用损失函数和光栅化调优策略，在不增加球谐度数的情况下处理视角相关颜色。训练速度约30倍快于RawNeRF（<=15分钟/场景），推理速度>=120fps。适用时还应保留论文给出的边界：需要RAW格式图像作为输入；对极端动态范围场景的处理仍有改进空间。
% METHOD: 2403.15704
\paragraph{Gaussian in the Wild（非受控照片集外观分解）（2024.03）}
Gaussian in the Wild~\cite{2403_15704} 是野外外观变化路线的代表性锚点。旅游照片集或互联网图像集合不是同一台相机、同一时间、同一曝光下拍摄的静态场景：天空、天气、阴影、高光会变，行人和车辆还会遮挡场景。原始 3DGS 如果把这些变化都写进固定 SH 颜色和几何，就会生成永久悬浮的“幽灵”高斯。GS-W 的改变是放弃用单一球谐颜色解释所有图像，为每个高斯拆出 intrinsic appearance feature 和 dynamic appearance feature。

具体做法是：每个高斯保留一个由材质和表面属性决定的内在外观特征，同时从参考图像中提取动态外观特征。网络先从参考图像得到一个 projection feature map，高斯中心通过已知相机投影到图像上，用双线性采样得到局部特征；但同一条射线上的高斯会采到相同位置，不能区分不同深度点受到的局部光照，因此 GS-W 额外提取 $K$ 张 feature maps，并给每个高斯学习 $K$ 个采样坐标，让它可以在这些特征图中自适应选择更有用的局部高频信息。最终，一个小型 MLP 融合内在特征和动态特征，输出该参考外观下的颜色。

瞬态物体通过 2D visibility map 处理。训练时，网络估计当前图像中哪些像素是静态场景可解释区域，哪些更可能来自行人、车辆等临时遮挡；光度损失按可见性加权，避免遮挡区域的错误梯度推动高斯生长或漂移。实现上，官方代码在 PhotoTourism 的 Brandenburg Gate、Trevi Fountain、Sacre Coeur 等场景上训练，保留 3DGS 的剪枝/增长机制以补足初始点云缺失，同时用外观网络处理不同照片的局部光照。相对 NeRF-W/Ha-NeRF 的全局 appearance embedding，GS-W 的关键是把外观条件下沉到每个高斯和局部图像特征；相对 Splatfacto-W 的工程化实现，它是更早的核心模型来源。
\textbf{锚点细化：} Gaussian in the Wild (GS-W)：针对互联网照片集合中光照变化和瞬态物体（行人、车辆）并存导致 3DGS 创建永久悬浮幽灵高斯的局限，用自适应外观特征（分解为固有外观 + 动态外观）+ 2D 可见性图替代球谐颜色，核心思想是为每个高斯引入外观特征，从参考图像通过 UNet-ResNet 提取投影特征和 $K$ 张附加特征图，每个高斯拥有 $K$ 个可学习采样坐标自适应决定在哪些位置采样动态外观特征，同时训练 UNet 生成 2D 可见性图加权光度损失掩蔽瞬态物体。
\textbf{方法信息：} 代码：\url{https://github.com/eastbeanzhang/gaussian-wild}；GitHub 关注度较高，约 200 stars；引用很多，Semantic Scholar 约 147 次。



% METHOD: 2501.10788
\paragraph{3D Consistent Appearance Features（外观变化解耦）（2025.01）}
Decoupling Appearance Variations with 3D Consistent Features in Gaussian Splatting~\cite{2501_10788} 也处理非受控采集下的外观变化，但强调 3D 一致特征。不同视角中的曝光、白平衡或阴影变化不应改变同一个 3D 点的身份；方法通过学习 3D 一致外观特征，把短期图像变化和长期场景属性分开。

DAVIGS 的关键选择是不在高斯属性里直接改颜色，而是在图像级做外观变换。基础 3DGS 先按普通方式渲染出图像；随后每个训练视图有一个全局 appearance embedding，用轻量 MLP 预测全局仿射颜色变换，解释相机 ISP、天气或曝光差异。对于局部光照变化，方法再引入多分辨率哈希网格作为 3D 一致空间特征：像素通过渲染深度反投影到 3D 坐标，在哈希网格中查询局部特征，再和全局 appearance embedding 一起生成像素级校正。这样，同一个 3D 位置在不同视角中查询到的局部特征一致，而每张图像仍可有自己的外观偏移。相对 GS-W，它更像一个可插拔后处理外观模块；相对 Luminance-GS++，它不仅做全局曲线，还显式利用 3D 一致局部特征处理局部光照变化。
\textbf{方法信息：} Huawei 参与了这篇论文；当前未确认完整公开代码。
\textbf{设计动机与边界：} 相对原始 3DGS 以独立高斯和图像重建误差为主的设定，Decoupling Appearance Variations with 3D Consistent Features in Gaussian Splatting 把问题落在以下可实现的改动上：DAVIGS方法的核心设计：(1)图像级解耦：不在高斯层面（如修改颜色/不透明度属性）进行外观变换，而是在图像渲染后进行变换，保持渲染管线不变。使用全局共享的外观嵌入（AE）向量编码每个视角的全局外观（如ISP差异、光照变化），通过轻量MLP学习仿射变换矩阵。(2)3D一致性局部特征：使用多分辨率哈希网格存储局部空间特征。对渲染图像的每个像素，通过反投影计算3D坐标，在哈希网格中查询周围体素并插值获取局部特征。结合全局AE和局部空间特征生成像素级变换。(3)即插即用：DAVIGS可以无缝集成到任何GS基线（3DGS、Mip-Splatting、2DGS等），优化完成后可丢弃外观模块，不影响渲染速度。其设计目的不是单独增加一个模块，而是用上述表示、监督或渲染约束消除原流程在目标任务中的具体失配。提出DAVIGS方法，在图像层面而非高斯层面解耦外观变化，结合3D一致性空间特征实现即插即用的外观建模。在训练时间、显存占用和渲染速度上开销最小，同时保持SOTA渲染质量，且兼容不同GS基线方法。适用时还应保留论文给出的边界：对极端局部光照变化（如聚光灯）的建模能力有限；哈希网格的分辨率影响局部特征的表达能力。
% METHOD: 2503.10143
\paragraph{GaussHDR（统一 3D/2D 局部色调映射）（2025.03）}
GaussHDR~\cite{2503_10143} 在 HDR 3DGS 中同时学习 3D 辐射表示和 2D 局部色调映射。仅用全局 tone mapping 很难处理局部高亮和暗部细节；只在 2D 后处理又会破坏多视角一致。GaussHDR 的改变是统一高斯级 HDR 颜色和视图级局部映射。

每个高斯不只携带 HDR radiance，还携带一个用于 tone mapping 的 context feature。3D local tone mapping 在高斯属性层面把 HDR 值映射成曝光相关 LDR；2D local tone mapping 则先把 context feature 随渲染投到图像平面，再在像素级处理局部亮度和对比度。二者各自得到一张 LDR 预测，训练时不简单平均，而是在损失层面同时计算 3D tone mapping 和 2D tone mapping 的重建误差，并用 uncertainty learning 自适应决定当前场景更信任哪一路。优化流程先用全局 tone mapper 稳定 HDR 辐射，再打开残差局部 tone mapper 和不确定性模块；输出既可以是 HDR 新视角，也可以是指定曝光下的 LDR 新视角。这样，3D 路径维持跨视角一致的 HDR 辐射，2D 路径补足局部 tone mapping 的表达力。相对 HDR-GS 的双动态范围高斯，GaussHDR 更强调“3D 辐射稳定”和“2D 局部相机响应”协同，而不是只学习一个全局 CRF。
\textbf{方法信息：} 代码：\url{https://github.com/liujf1226/gausshdr}。
\textbf{设计动机与边界：} 相对原始 3DGS 以独立高斯和图像重建误差为主的设定，GaussHDR 把问题落在以下可实现的改动上：GaussHDR基于3DGS，将HDR辐照度替换LDR颜色属性。核心创新：(1) 残差局部色调映射器同时支持3D和2D色调映射，接受额外的上下文特征输入以实现局部化处理；(2) 双LDR渲染结果结合策略，在损失级别融合3D色调映射（LDR拟合好但HDR不稳定）和2D色调映射（HDR稳定但LDR拟合差）的互补优势；(3) 不确定性学习模块，根据不同场景中双结果的平衡差异进行自适应调制。在球谐函数表示颜色基础上，为RGB通道独立建模色调映射器。其设计目的不是单独增加一个模块，而是用上述表示、监督或渲染约束消除原流程在目标任务中的具体失配。首次统一3D和2D局部色调映射用于HDR NVS，提出残差局部色调映射器和不确定性自适应调制策略，在合成和真实场景中显著超越SOTA方法。适用时还应保留论文给出的边界：局部色调映射器增加了参数量；不确定性学习的额外开销；对曝光时间的准确标注有依赖。
% METHOD: 2602.18322
\paragraph{Luminance-GS++（视角自适应曲线校正）（2026.02）}
Unifying Color and Lightness Correction with View-Adaptive Curve Adjustment~\cite{2602_18322} 处理低光、过曝和复杂亮度/色度变化。它保留显式 3DGS 表示，不引入完整材质分解，而是在颜色输出端增加视角自适应校正。

全局模块学习随视角变化的亮度曲线调整函数，可用分段线性或非线性曲线把每个视角的整体亮度校正到一致范围。局部模块在全局调整后加入轻量像素级残差网络，补偿局部光照和色彩偏差。无监督损失同时约束亮度一致、光度重建和几何一致。相对 GS-W 的每高斯动态特征，它更像把图像处理曲线嵌入 3DGS 训练；相对 HDR-GS，它不追求恢复物理 HDR radiance，而是让输入视角在低光、过曝和白平衡不一致时仍能形成统一的 NVS 训练目标。
\textbf{实现细节：} Unifying Color and Lightness Correction（2026.02） （View-Adaptive Curve Adjustment）解决不同训练图像之间曝光和白平衡不一致的问题。标准 3DGS 假设所有图像在统一颜色空间中，但实际上每张照片的相机 ISP 处理链不同。该方法为每张训练图像学习一组色调曲线参数，通过视角自适应的曲线调整将不同曝光和白平衡的图像映射到统一颜色空间，使 3DGS 的光度损失在一致的色彩空间中计算。
\textbf{锚点细化：} Unifying Color and Lightness Correction：（View-Adaptive Curve Adjustment）解决不同训练图像之间曝光和白平衡不一致的问题。标准 3DGS 假设所有图像在统一颜色空间中，但实际上每张照片的相机 ISP 处理链不同。该方法为每张训练图像学习一组色调曲线参数，通过视角自适应的曲线调整将不同曝光和白平衡的图像映射到统一颜色空间，使 3DGS 的光度损失在一致的色彩空间中计算。
\textbf{设计动机与边界：} 相对原始 3DGS 以独立高斯和图像重建误差为主的设定，Unifying Color and Lightness Correction with View-Adaptive Curve Adjustment for Robust 3D  把问题落在以下可实现的改动上：框架基于3DGS扩展，保持显式高斯表示。两个核心模块：(1)全局视角自适应亮度调整：设计可学习的曲线调整函数，根据视角参数自适应调整亮度。使用分段线性或非线性曲线映射，对每个视角的全局亮度进行校正。曲线参数与高斯参数联合优化；(2)局部像素级残差细化：在全局调整基础上，添加轻量级像素级残差网络，捕获局部光照变化和色彩偏差。无监督损失设计：联合亮度一致性损失（确保多视角亮度一致）、光度一致性损失（3DGS渲染损失）和几何一致性损失。无需GT光照条件监督。其设计目的不是单独增加一个模块，而是用上述表示、监督或渲染约束消除原流程在目标任务中的具体失配。提出Luminance-GS++，结合全局视角自适应亮度调整和局部像素级残差细化实现精确颜色校正。设计无监督目标联合强制亮度校正和多视角几何/光度一致性。保持显式3DGS公式，在低光、过曝光和复杂亮度/色度变化场景下达到SOTA性能。适用时还应保留论文给出的边界：曲线参数的容量有限；极端光照变化下的校正能力可能不足；增加的训练开销。
\textbf{方法信息：} Google DeepMind 参与了这篇论文；作者关联代码：\url{https://github.com/cuiziteng/Luminance-GS}；该仓库按年龄归一后 stars 表现突出。



\section{基于物理的渲染与光照}

原始 3DGS 把外观直接拟合到颜色和球谐系数中，速度很快，但光照、材质和多次反弹都被混在一起。基于物理的路线要做的改变，是把“看起来像”的辐射场拆成可解释的反照率、粗糙度、发射、阴影和全局光照，并让高斯能够进入路径追踪、可重光照或编辑流程。这里的关键不只是画面更真实，而是换光、移物体、改材质时结果仍然合理。



% METHOD: 2601.23065
\paragraph{EAG-PT（发射感知高斯与路径追踪）（2026.01）}
EAG-PT~\cite{2601_23065} 把 2D 高斯重新表述为可路径追踪的几何代理。普通 3DGS/2DGS 很快，但不支持物理正确多次反弹；纯 mesh 路线又要显式重建。EAG-PT 的改变是用 2DGS 表示室内场景，并分离发射和非发射组件。

它先把 2D 高斯写成可求交的椭圆面元：射线与高斯相交后按响应函数和透明度做 alpha 累积，得到法线、深度、辐射、发射项和反照率等宏观量。重建阶段分两步。Stage 0 在无反弹 tracing 下恢复 capture-time radiance，并用 2D emission mask 明确区分灯具等发射区域和非发射几何；这里的训练损失包含图像重建、法线/深度稳定项和发射约束。Stage 1 在固定几何上引入一次反弹 differentiable rendering，通过半球方向采样和发射器辐射估计入射光，用场景 radiance cache 反推出非发射区域的 diffuse albedo。编辑或换灯后，最终图像不再靠原来的 radiance cache，而是用多次反弹 path tracing 重新计算间接光；为了交互浏览，还可把 path-traced 结果重新 bake 回高斯辐射中。相对 GI-GS 的间接光补偿，EAG-PT 更彻底地把高斯作为 traceable primitive 接入路径追踪，但当前材料主要是漫反射，镜面和折射仍是后续方向。
\textbf{实现细节：} EAG-PT（2026.01） （Emission-Aware Gaussians and Path Tracing）解决的核心问题是室内场景的全局光照。标准 3DGS 的 $\alpha$-混合 假设光线只沿直线传播，无法处理间接光照（如灯光照亮墙壁后反射到地板）。EAG-PT 采用两阶段管线。Stage 0（辐射场景重建）：使用 2D 高斯作为场景表示，每个 2D 高斯在交点处的响应为 $g(\mathbf{x}) = \exp\left(-\frac{1}{2}\left(\frac{(\mathbf{x}-\mathbf{p}_u)^{}2}{s_u^{}2} + \frac{(\mathbf{x}-\mathbf{p}_v)^{}2}{s_v^{}2}\right)\right)$，其中 $s_{u,v}$ 为各向异性尺度参数。通过可微 2D 高斯光线追踪渲染辐射图 $L_c^{}0 = L_c(\mathrm{PQ}(L_o^{}0), \mathrm{PQ}(L_o^{\mathrm{gt}}))$，其中 $\mathrm{PQ}(\cdot)$ 是 SMPTE ST-2084 感知量化曲线。几何损失用单目法线图监督：$L_n = \|1 - \hat{N} \cdot N_{\mathrm{mono}}\|_1$。发光损失用 2D 发光遮罩 $M$ 监督：$L_e = \|E - M\|_1$。
\textbf{技术实现补充：} EAG-PT（2026.01） （Emission-Aware Gaussians and Path Tracing）解决的核心问题是室内场景的全局光照。标准 3DGS 的 $\alpha$-混合 假设光线只沿直线传播，无法处理间接光照（如灯光照亮墙壁后反射到地板）。EAG-PT 采用两阶段管线。Stage 0（辐射场景重建）：使用 2D 高斯作为场景表示，每个 2D 高斯在交点处的响应为 $g(\mathbf{x}) = \exp\left(-\frac{1}{2}\left(\frac{(\mathbf{x}-\mathbf{p}_u)^2}{s_u^2} + \frac{(\mathbf{x}-\mathbf{p}_v)^2}{s_v^2}\right)\right)$，其中 $s_{u,v}$ 为各向异性尺度参数。通过可微 2D 高斯光线追踪渲染辐射图 $L_c^0 = L_c(\mathrm{PQ}(L_o^0), \mathrm{PQ}(L_o^{\mathrm{gt}}))$，其中 $\mathrm{PQ}(\cdot)$ 是 SMPTE ST-2084 感知量化曲线。几何损失用单目法线图监督：$L_n = \|1 - \hat{N} \cdot N_{\mathrm{mono}}\|_1$。发光损失用 2D 发光遮罩 $M$ 监督：$L_e = \|E - M\|_1$。Stage 1（漫反射材质恢复）：单次反弹近似间接光照——将 Stage 0 重建的辐射场作为辐射缓存（radiance cache），通过蒙特卡洛积分 $L_r^1(\mathbf{d}_o) \approx \frac{1}{n_{\mathrm{spp}}} \sum^{n_{\mathrm{spp}}} R(\mathbf{d}_i) f(\mathbf{d}_i, \mathbf{d}_o) (\mathbf{d}_i \cdot \mathbf{N}) / p(\mathbf{d}_i)$ 估计漫反射反照率 $\rho$。假设发射表面不反射（$E(\mathbf{d}_o) > \epsilon_E=0.1$ 时仅发光），实现了发光和反射的分离。场景编辑后使用多反弹路径追踪产生物理正确渲染，并通过光照烘焙将路径追踪结果回写到高斯的辐射值以实现实时预览。
\textbf{方法信息：} 代码：\url{https://github.com/internrobotics/eag-pt}。
\textbf{设计动机与边界：} 相对原始 3DGS 以独立高斯和图像重建误差为主的设定，EAG-PT 把问题落在以下可实现的改动上：EAG-PT基于2D高斯表示，包含四个紧密耦合的组件：(1) 发射感知场景分解：使用2D发射掩码显式分离发射光源和非发射几何，解决以往辐射场景模型无法支持编辑的问题。(2) 辐射度和材质属性的逆向恢复：通过不同iable渲染恢复发射器辐射度和非发射区域反照率。(3) 单次反弹逆向优化：阶段0使用纯追踪（无光反弹）进行辐射场景重建；阶段1引入渲染方程，通过单次反弹的differentiable渲染估计漫反射材质属性。(4) 多次反弹路径追踪：用于正向渲染，结合光烘焙技术实现物理正确的全局光照。整体流程：多视图图像→2DGS辐射场景→材质恢复→场景编辑→路径追踪渲染。2D高斯被重新表述为可追踪基元，支持射线-高斯交点的解析计算。其设计目的不是单独增加一个模块，而是用上述表示、监督或渲染约束消除原流程在目标任务中的具体失配。提出EAG-PT，基于统一2D高斯表示的物理室内场景重建与编辑方法。核心创新：(1) 用2D高斯作为传输友好的几何代理，避免显式网格重建；(2) 重建时显式分离发射和非发射组件支持编辑；(3) 解耦重建（高效单次反弹优化）和最终渲染（高质量多次反弹路径追踪）。适用时还应保留论文给出的边界：当前仅支持漫反射材质；对高光和镜面反射的处理有限；依赖密集多视图输入和线性色彩空间；路径追踪渲染开销较高；初始化质量影响最终结果。

\section{可变形与可重光照资产}

可变形资产把逆渲染问题再推进一步：对象不再只是静态场景，而是之后还要换姿态、换材质、换光照或插入其他环境。普通 3DGS 把外观烘焙在高斯上，形变后容易出现阴影和高光不跟随表面的情况；本节方法把材质、光照和几何绑定关系拆开，使高斯更接近可编辑资产而不是一次性重建结果。
% METHOD: 2311.16043
\paragraph{Relightable 3D Gaussians（逆渲染、重光照与材质分解）（2023.11）}
Relightable 3D Gaussians~\cite{2311_16043} 关注的问题是：把 3DGS 扩展到逆渲染、材质分解和重光照 Relightable 3D Gaussians 是宽召回补全中发现的论文，主要根据标题、摘要和 arXiv 类别判定其与 3DGS 的关系。
相对原始 3DGS，Relightable 3D Gaussians 的变化不是只替换损失函数，而是把表示、初始化、约束或网络预测路径重新组织到任务流程中：该工作被召回是因为论文文本包含 Gaussian/3D Gaussian 相关表述；当前分类先确定其任务域和与 3DGS 的相关性，后续证据流水线再补代码、企业和引用信息。
\textbf{实现细节：} 具体实现可以按输入、核心模块和渲染/优化路径来读。输入端首先取得任务所需的图像、相机、文本、时间序列、语义或几何线索，并把这些线索转换成可以约束高斯中心、尺度、旋转、不透明度、颜色或附加特征的条件。核心结构中，方法使用的主要机制是：从摘要层面看，Relightable 3D Gaussians 围绕“把 3DGS 扩展到逆渲染、材质分解和重光照”展开。若其正文或代码后续可读，应进一步核对它具体改动的是表示、初始化、渲染、优化、损失、动态建模、压缩或应用接口中的哪一环。 训练或推理时，这些模块通常先给出初始化或中间特征，再通过可微 Gaussian rasterization、重建损失、感知/语义约束、几何正则或任务损失回传到高斯参数；如果方法包含前馈网络，则网络输出的是高斯属性、变形场、压缩码、语义特征或控制变量，而不是只输出最终 2D 图像。渲染路径仍以显式高斯投影、深度排序和 alpha 合成为基础，只是在采样、属性预测、局部编辑、动态变形、压缩编码或下游任务接口处加入额外约束。这样写入本章时，它与前后方法的关系是：相对原始 3DGS 和本章已有锚点方法，它可以理解为 3DGS extends\_or\_uses 宽召回分类阶段的关系判断，后续需要按论文方法部分复核。
\textbf{证据与边界：} 本轮召回分类尚未逐项核对实验表格；后续证据补全阶段需要读取论文源码/PDF或项目页确认数据集、指标和对比方法。 该条来自宽召回批次；源码下载受 arXiv e-print 网络限制，部分论文暂时只能依据标题、摘要和类别分类，需要在入可信书前做人工/LLM精读复核。 本书只采用本地 paper pack、数据库分类、源码/引用登记和论文摘要中可核验的信息；没有本地证据支撑的速度、排行榜或开源完整性不在正文中强化。\textbf{方法信息：} Huawei 参与；代码：\url{https://github.com/NJU-3DV/Relightable3DGaussian}；引用信号：169 citations (5.4/month, domain-relative cutoff 2.2/month, semantic\_scholar)；代码关注：nju-3dv/relightable3dgaussian has 677 stars (21.5/month, domain-relative cutoff 12.8/month)。
\textbf{设计动机与边界：} 相对原始 3DGS 以独立高斯和图像重建误差为主的设定，Relightable 3D Gaussians 把问题落在以下可实现的改动上：从摘要层面看，Relightable 3D Gaussians 围绕“把 3DGS 扩展到逆渲染、材质分解和重光照”展开。若其正文或代码后续可读，应进一步核对它具体改动的是表示、初始化、渲染、优化、损失、动态建模、压缩或应用接口中的哪一环。其设计目的不是单独增加一个模块，而是用上述表示、监督或渲染约束消除原流程在目标任务中的具体失配。该工作被召回是因为论文文本包含 Gaussian/3D Gaussian 相关表述；当前分类先确定其任务域和与 3DGS 的相关性，后续证据流水线再补代码、企业和引用信息。适用时还应保留论文给出的边界：该条来自宽召回批次；源码下载受 arXiv e-print 网络限制，部分论文暂时只能依据标题、摘要和类别分类，需要在入可信书前做人工/LLM精读复核。
% METHOD: 2312.05133
\paragraph{GIR（逆渲染、重光照与材质分解）（2023.12）}
GIR~\cite{2312_05133} 关注的问题是：把 3DGS 扩展到逆渲染、材质分解和重光照 GIR 是宽召回补全中发现的论文，主要根据标题、摘要和 arXiv 类别判定其与 3DGS 的关系。
相对原始 3DGS，GIR 的变化不是只替换损失函数，而是把表示、初始化、约束或网络预测路径重新组织到任务流程中：该工作被召回是因为论文文本包含 Gaussian/3D Gaussian 相关表述；当前分类先确定其任务域和与 3DGS 的相关性，后续证据流水线再补代码、企业和引用信息。
\textbf{实现细节：} 具体实现可以按输入、核心模块和渲染/优化路径来读。输入端首先取得任务所需的图像、相机、文本、时间序列、语义或几何线索，并把这些线索转换成可以约束高斯中心、尺度、旋转、不透明度、颜色或附加特征的条件。核心结构中，方法使用的主要机制是：从摘要层面看，GIR 围绕“把 3DGS 扩展到逆渲染、材质分解和重光照”展开。若其正文或代码后续可读，应进一步核对它具体改动的是表示、初始化、渲染、优化、损失、动态建模、压缩或应用接口中的哪一环。 训练或推理时，这些模块通常先给出初始化或中间特征，再通过可微 Gaussian rasterization、重建损失、感知/语义约束、几何正则或任务损失回传到高斯参数；如果方法包含前馈网络，则网络输出的是高斯属性、变形场、压缩码、语义特征或控制变量，而不是只输出最终 2D 图像。渲染路径仍以显式高斯投影、深度排序和 alpha 合成为基础，只是在采样、属性预测、局部编辑、动态变形、压缩编码或下游任务接口处加入额外约束。这样写入本章时，它与前后方法的关系是：相对原始 3DGS 和本章已有锚点方法，它可以理解为 3DGS extends\_or\_uses 宽召回分类阶段的关系判断，后续需要按论文方法部分复核。
\textbf{证据与边界：} 本轮召回分类尚未逐项核对实验表格；后续证据补全阶段需要读取论文源码/PDF或项目页确认数据集、指标和对比方法。 该条来自宽召回批次；源码下载受 arXiv e-print 网络限制，部分论文暂时只能依据标题、摘要和类别分类，需要在入可信书前做人工/LLM精读复核。 本书只采用本地 paper pack、数据库分类、源码/引用登记和论文摘要中可核验的信息；没有本地证据支撑的速度、排行榜或开源完整性不在正文中强化。\textbf{方法信息：} 代码：\url{https://github.com/guduxiaolang/gir}；引用信号：105 citations (3.4/month, domain-relative cutoff 2.2/month, semantic\_scholar)。
\textbf{设计动机与边界：} 相对原始 3DGS 以独立高斯和图像重建误差为主的设定，GIR 把问题落在以下可实现的改动上：从摘要层面看，GIR 围绕“把 3DGS 扩展到逆渲染、材质分解和重光照”展开。若其正文或代码后续可读，应进一步核对它具体改动的是表示、初始化、渲染、优化、损失、动态建模、压缩或应用接口中的哪一环。其设计目的不是单独增加一个模块，而是用上述表示、监督或渲染约束消除原流程在目标任务中的具体失配。该工作被召回是因为论文文本包含 Gaussian/3D Gaussian 相关表述；当前分类先确定其任务域和与 3DGS 的相关性，后续证据流水线再补代码、企业和引用信息。适用时还应保留论文给出的边界：该条来自宽召回批次；源码下载受 arXiv e-print 网络限制，部分论文暂时只能依据标题、摘要和类别分类，需要在入可信书前做人工/LLM精读复核。
% METHOD: 2402.10483
\paragraph{GaussianHair（逆渲染、重光照与材质分解）（2024.02）}
GaussianHair~\cite{2402_10483} 关注的问题是：把 3DGS 扩展到逆渲染、材质分解和重光照 GaussianHair 是宽召回补全中发现的论文，主要根据标题、摘要和 arXiv 类别判定其与 3DGS 的关系。
相对原始 3DGS，GaussianHair 的变化不是只替换损失函数，而是把表示、初始化、约束或网络预测路径重新组织到任务流程中：该工作被召回是因为论文文本包含 Gaussian/3D Gaussian 相关表述；当前分类先确定其任务域和与 3DGS 的相关性，后续证据流水线再补代码、企业和引用信息。
\textbf{实现细节：} 具体实现可以按输入、核心模块和渲染/优化路径来读。输入端首先取得任务所需的图像、相机、文本、时间序列、语义或几何线索，并把这些线索转换成可以约束高斯中心、尺度、旋转、不透明度、颜色或附加特征的条件。核心结构中，方法使用的主要机制是：从摘要层面看，GaussianHair 围绕“把 3DGS 扩展到逆渲染、材质分解和重光照”展开。若其正文或代码后续可读，应进一步核对它具体改动的是表示、初始化、渲染、优化、损失、动态建模、压缩或应用接口中的哪一环。 训练或推理时，这些模块通常先给出初始化或中间特征，再通过可微 Gaussian rasterization、重建损失、感知/语义约束、几何正则或任务损失回传到高斯参数；如果方法包含前馈网络，则网络输出的是高斯属性、变形场、压缩码、语义特征或控制变量，而不是只输出最终 2D 图像。渲染路径仍以显式高斯投影、深度排序和 alpha 合成为基础，只是在采样、属性预测、局部编辑、动态变形、压缩编码或下游任务接口处加入额外约束。这样写入本章时，它与前后方法的关系是：相对原始 3DGS 和本章已有锚点方法，它可以理解为 3DGS extends\_or\_uses 宽召回分类阶段的关系判断，后续需要按论文方法部分复核。
\textbf{证据与边界：} 本轮召回分类尚未逐项核对实验表格；后续证据补全阶段需要读取论文源码/PDF或项目页确认数据集、指标和对比方法。 该条来自宽召回批次；源码下载受 arXiv e-print 网络限制，部分论文暂时只能依据标题、摘要和类别分类，需要在入可信书前做人工/LLM精读复核。 本书只采用本地 paper pack、数据库分类、源码/引用登记和论文摘要中可核验的信息；没有本地证据支撑的速度、排行榜或开源完整性不在正文中强化。\textbf{方法信息：} 公开链接：\url{https://github.com/HaiminLuo/GaussianHair}，完整性按本地证据保守处理；引用信号：62 citations (2.2/month, domain-relative cutoff 2.2/month, semantic\_scholar)。
\textbf{设计动机与边界：} 相对原始 3DGS 以独立高斯和图像重建误差为主的设定，GaussianHair 把问题落在以下可实现的改动上：从摘要层面看，GaussianHair 围绕“把 3DGS 扩展到逆渲染、材质分解和重光照”展开。若其正文或代码后续可读，应进一步核对它具体改动的是表示、初始化、渲染、优化、损失、动态建模、压缩或应用接口中的哪一环。其设计目的不是单独增加一个模块，而是用上述表示、监督或渲染约束消除原流程在目标任务中的具体失配。该工作被召回是因为论文文本包含 Gaussian/3D Gaussian 相关表述；当前分类先确定其任务域和与 3DGS 的相关性，后续证据流水线再补代码、企业和引用信息。适用时还应保留论文给出的边界：该条来自宽召回批次；源码下载受 arXiv e-print 网络限制，部分论文暂时只能依据标题、摘要和类别分类，需要在入可信书前做人工/LLM精读复核。
% METHOD: 2606.29379
\paragraph{DR-GS（可变形、可材质编辑、可重光照 2D 高斯）（2026.06）}
DR-GS~\cite{2606_29379} 处理动态资产中的材质和光照一致性。它继承 2DGS/高斯表面逆渲染的可微表面表达，同时接上 PhysGaussian、GSP 这类可变形 GS 的驱动路线。很多可变形 GS 方法先重建一份静态高斯，再用粒子或 mesh 驱动高斯运动；这样能让资产动起来，但颜色里已经烘焙了原始光照，物体弯曲、换姿态或换环境光后，高光和阴影不会按物理规律重新分布。DR-GS 的改变是把 2D Gaussians 用作可编辑表面资产，同时显式分离几何、材质和照明。

流程分为三个层次。第一层是静态重建：从多视图图像得到 2D Gaussian 表面，并提取 TSDF mesh 作为后续材质和光线追踪的几何载体。第二层是逆渲染：不再让高斯直接存最终 RGB，而是估计 diffuse/specular 材质、环境光和几何参数；渲染时求解带低采样 Monte Carlo 的 rendering equation，并用 multiple importance sampling 组合 cosine-weighted、GGX 和 environment-light 采样，再用 cross-bilateral denoising 降噪。第三层是动态驱动：模拟和渲染解耦，粒子系统或 mesh deformation 先给出控制点/网格形变，再用 generalized moving least squares 把形变插值到 2D Gaussians，使材质参数随表面走，而不是随屏幕颜色走。

它与 LumiMotion 的差别在于，LumiMotion 利用真实场景动态提供额外观察来做逆渲染，DR-GS 则面向内容生产中的“重建后还能编辑和动画化”的资产。它与 EAG-PT 的差别在于，DR-GS 为了连续形变的效率，把 decoupled material parameters 存在 TSDF mesh vertices 上，并在动态阶段用 mesh-based ray tracing 替代直接对 2D Gaussians 做重光线追踪。论文实验使用 GlossySynthetic、TensoIR、Sketchfab/Mixamo 角色等，评价静态重建、relighting、physical plausibility、temporal coherence 和 deformation consistency。当前本地证据显示 Huawei 参与，尚未记录到完整公开代码。
\textbf{实现细节：} DR-GS（2026.06） 延续 Gaussian Splatting 在"Deformable and relightable Gaussian assets"上的基本设定，主要改变是：物理逆渲染、relighting 与 deformation-aware manipulation 的统一 2D Gaussian 框架。3DGS 的显式高斯表示和可微渲染能力保持为基础。其设计动机是：解决既有 GS 方法把照明烘焙进纹理、难以在动态形变和光照变化下保持物理一致的问题。具体而言，以 2D Gaussians 为显式表示，联合估计材质、光照和可形变结构，使重建后仍可进行材质编辑、动态变形和重光照。适用边界是：物理逆渲染欠定，复杂 BRDF、遮挡和非刚性形变可能带来估计不稳定。
\textbf{实现补充：} 实现上还要核对输入、输出、参数、特征、初始化、训练损失、正则、渲染或推理接口，以及优化中的剪枝、分裂、采样、编码或解码步骤；这些抓手决定该方法到底改了 3DGS 管线的哪一环。
\textbf{技术实现补充：} DR-GS。DR-GS（2026.06） 延续 Gaussian Splatting 在"Deformable and relightable Gaussian assets"上的基本设定，主要改变是：物理逆渲染、relighting 与 deformation-aware manipulation 的统一 2D Gaussian 框架。3DGS 的显式高斯表示和可微渲染能力保持为基础。其设计动机是：解决既有 GS 方法把照明烘焙进纹理、难以在动态形变和光照变化下保持物理一致的问题。具体而言，以 2D Gaussians 为显式表示，联合估计材质、光照和可形变结构，使重建后仍可进行材质编辑、动态变形和重光照。适用边界是：物理逆渲染欠定，复杂 BRDF、遮挡和非刚性形变可能带来估计不稳定。
\textbf{方法信息：} Huawei 参与了这篇论文。

\begin{takeaway}
本章的方法可以按解耦对象分层理解。DeferredGS 奠定了“先 G-buffer 后着色”的高斯逆渲染思路；Ref-Gaussian、RGS-DR、MaterialRefGS 和 Spec-Gloss Surfels 沿着这条路强化反射、BRDF 和材质一致性；Mirror-3DGS、Flash-Splat、Latent-SpecGS 和 Augmented Radiance Field 则专门处理镜面、玻璃和尖锐高光。重光照方向从 LumiGauss/BiGS/RNG 的光照传输发展到 RelitLRM/F-RNG 的前馈可重光照资产。HDR 和野外外观方向则由 GS-W、HDR-GS 与 LE3D 代表，分别处理采集条件变化、多曝光动态范围和 RAW 低光线性空间。整体看，3DGS 从“把光照烘焙进颜色的快速渲染器”逐步变成“能拆材质、拆光照、拆相机响应，并接入物理渲染的场景表示”。
\end{takeaway}
\textbf{设计动机与边界：} 相对原始 3DGS 以独立高斯和图像重建误差为主的设定，DR-GS 把问题落在以下可实现的改动上：以 2D Gaussians 为显式表示，联合估计材质、光照和可形变结构，使重建后仍可进行材质编辑、动态变形和重光照。其设计目的不是单独增加一个模块，而是用上述表示、监督或渲染约束消除原流程在目标任务中的具体失配。解决既有 GS 方法把照明烘焙进纹理、难以在动态形变和光照变化下保持物理一致的问题。适用时还应保留论文给出的边界：物理逆渲染欠定，复杂 BRDF、遮挡和非刚性形变可能带来估计不稳定。
